\chapter{Hidden extensions}
\label{ch:hidden}

In this chapter, we will discuss hidden extensions
in the $E_\infty$-page of the Adams spectral sequence.
We methodically explore hidden extensions by
$\tau$, $2$, $\eta$, and $\nu$, and we study other miscellaneous
hidden extensions that are relevant for specific purposes.
\revv{
For easy reference, the lemmas in this chapter are labelled with degrees that match the degrees given in the tables.}

\section{Hidden $\tau$ extensions}
\label{sctn:tau-extn}

In order to study hidden $\tau$ extensions,
we will use the long exact sequence
\begin{equation}
\label{eq:LES-Ctau}
\xymatrix@1{
\cdots \ar[r] & \pi_{p,q+1} \ar[r]^\tau & \pi_{p,q} \ar[r] & 
\pi_{p,q} C\tau \ar[r] & 
\pi_{p-1,q+1} \ar[r]^\tau & \pi_{p-1,q} \ar[r] & \cdots
}
\end{equation}
extensively.
This sequence governs hidden $\tau$ extensions in the following sense.
An element $\alpha$ in $\pi_{p,q}$ is divisible by $\tau$
if and only if it maps to zero in $\pi_{p,q} C\tau$,
and an element $\alpha$ in $\pi_{p-1,q+1}$ supports a $\tau$ extension
if and only if it is not in the image of $\pi_{p,q} C\tau$.
Therefore, we need to study the maps
$\pi_{*,*} \map \pi_{*,*} C\tau$ and
$\pi_{*,*} C\tau \map \pi_{*-1,*+1}$
induced by inclusion of the bottom cell into $C\tau$
and by projection from $C\tau$ to the top cell.

The $E_\infty$-pages of the Adams spectral sequences for $S^{0,0}$ and
$C\tau$ give associated graded objects for the homotopy groups that 
are the sources and targets of these maps.  Naturality of the
Adams spectral sequence induces maps on associated graded objects.

These maps on associated graded objects often detect the values
of the maps on homotopy groups.  For example, the element
$h_0$ in the Adams spectral sequence for the sphere is mapped to the
element $h_0$ in the Adams spectral sequence for $C\tau$.
In homotopy groups, this means that inclusion of the bottom cell into
$C\tau$ takes the element $2$ in $\pi_{0,0}$ to the
element $2$ in $\pi_{0,0} C\tau$.

On the other side, the element $\ol{h_1^4}$ in the 
Adams spectral sequence for $C\tau$ is mapped to the
element $h_1^4$ in Adams spectral sequence for the sphere.
In homotopy groups, this means that projection from $C\tau$ to
the top cell takes the element $\{ \ol{h_1^4} \}$ in $\pi_{5,3}C\tau$
to the element $\eta^4$ in $\pi_{4,4}$.

However, some values of the maps on homotopy groups can be hidden
in the map of associated graded objects.  This situation is rare in
low stems but becomes more and more common in higher stems.
The first such example occurs in the 30-stem.
The element $\D h_2^2 $ is a permanent cycle in the Adams spectral
sequence for $C\tau$, so $\{\D h_2^2 \}$ is an element in $\pi_{30,16} C\tau$.
Now $\D h_2^2$ maps to zero in the $E_\infty$-page of the Adams spectral
sequence for the sphere, but $\{\D h_2^2 \}$ does not map to zero
in $\pi_{29,17}$.  In fact $\{\D h_2^2 \}$ maps to $\eta \kappa^2$,
which is detected by $h_1 d_0^2$.
This demonstrates that projection from $C\tau$ to the top cell
has a hidden value.

We refer the reader to Section \ref{sctn:assoc-graded}
for a precise discussion of these issues.

\begin{thm}
\label{thm:hid-incl-proj}
\mbox{}
\begin{enumerate}
\item
Through the 90-stem,
Table \ref{tab:hid-incl} lists all
hidden values of inclusion of the bottom cell into $C\tau$,
except that:
\rev{
\begin{enumerate}
\item
If $h_1 f_2$ does not survive but $\tau h_1 f_2$ does survive,
then $\tau h_1 f_2$ maps to $h_0^3 c_3$.
\item
If $h_1^2 f_2$ does not survive, then $\tau h_1^2 f_2$
maps to $\ol{\tau h_1^3 h_4 Q_3}$ or
$\D^2 e_1 + \ol{\tau \D h_2 e_1 g}$.
\item
$\tau h_1 x_{85,6}$
maps to $\ol{\tau h_1^3 h_4 Q_3}$ or
$\D^2 e_1 + \ol{\tau \D h_2 e_1 g}$.
\end{enumerate}
}
\item
Through the 90-stem, Table \ref{tab:hid-proj} lists all
hidden values of projection from $C\tau$ to the top cell,
except that:
\rev{
\begin{enumerate}
\item
If $h_1 f_2$ does not survive, then $h_1 f_2$ maps to $h_1^2 h_4 Q_3$ 
or $\D h_1 j_1$.
\item
$x_{85,6}$ maps to $h_1^2 h_4 Q_3$ or $\D h_1 j_1$.
\item
If $h_1^2 f_2$ does not survive, then $h_1^2 f_2$ maps to
$h_1^3 h_4 Q_3$ or $\tau M h_0 g^2$.
\item
$h_1 x_{85,6}$ maps to $h_1^3 h_4 Q_3$ or $\tau M h_0 g^2$.
\item
If $h_1^2 f_2$ survives, then 
$\ol{\tau h_1^3 h_4 Q_3}$ or $\D^2 e_1 + \ol{\tau \D h_2 e_1 g}$
maps to $M \D h_1^2 d_0$.
\end{enumerate}
}
\end{enumerate}
\end{thm}

\begin{proof}
The values of inclusion of the bottom cell and projection
to the top cell are almost entirely determined by inspection
of Adams $E_\infty$-pages.
Taking into account the multiplicative structure,
there are no other combinatorial possibilities.
For example, consider the exact sequence
\[
\pi_{30,16} \map \pi_{30,16} C\tau \map \pi_{29,17}.
\]
In the Adams $E_\infty$-page for $C\tau$, 
$h_4^2$ and $\D h_2^2$
are the only two elements in the 30-stem with weight 16.
In the Adams $E_\infty$-page for the sphere,
$h_4^2$ is the only element in the 30-stem with weight 16,
and $h_1 d_0^2$ is the only element in the 29-stem with weight 17.
The only possibility is that $h_4^2$ maps to $h_4^2$ under
inclusion of the bottom cell, and
$\D h_2^2$ maps to $h_1 d_0^2$ under projection to the top cell.

One case, given below in Lemma \ref{lem:t-d1e1}, 
requires a more complicated argument.
\end{proof}

\begin{remark}
Through the 90-stem, inclusion of the bottom cell into $C\tau$
has only one hidden value with target indeterminacy.
Namely, $h_2 c_1 A'$ is the hidden value of $\ol{h_1 g B_7}$,
with target indeterminacy generated by $\D j_1$.
Through the 90-stem, projection from $C\tau$ to the top cell
has no hidden values with target indeterminacy.
\end{remark}

\begin{remark}
Through the 90-stem,
inclusion of the bottom cell into $C\tau$ has no crossing values.
On the other hand, projection from $C\tau$ to the top cell
does have crossing values in this range.  These occurrences
are described in the fourth column of 
Table \ref{tab:hid-proj}.
Each can be verified by direct inspection.
\end{remark}

\begin{thm}
\label{thm:tau-extn}
Through the 90-stem,
Table \ref{tab:tau-extn}
lists all hidden $\tau$ extensions in $\C$-motivic stable homotopy groups,
except that:
\rev{
\begin{enumerate}
\item
if $M \D h_1 d_0$ is not hit by a differential, then
there is a hidden $\tau$ extension from $\D h_1 j_1$ to
$M \D h_1 d_0$.
\item
if $M \D h_1^2 d_0$ is not hit by a differential, then
there is a hidden $\tau$ extension from $\tau M h_0 g^2$ to
$M \D h_1^2 d_0$.
\end{enumerate}
}
In this range, the only crossing extension is:
\begin{enumerate}
\item
the hidden $\tau$ extension from $h_1^2 h_6 c_0$ to $h_0 h_4 D_2$,
and the not hidden $\tau$ extension on $\tau h_2^2 Q_3$.
\end{enumerate}
\end{thm}

\begin{proof}
Almost all of these hidden $\tau$ extensions follow immediately
from the values of the maps in the long exact sequence
(\ref{eq:LES-Ctau})
given in Tables \ref{tab:hid-incl} and \ref{tab:hid-proj}.

For example, 
consider the element $P d_0$
in the Adams $E_\infty$-page for the sphere, which 
belongs to the 22-stem with weight 12. 
Now $\pi_{22,12} C\tau$ is zero because there are 
no elements in that degree in the Adams $E_\infty$-page
for $C\tau$, so
inclusion of the bottom cell takes $\{ P d_0 \}$ to zero.
Therefore, $\{ P d_0 \}$ must be in the image of multiplication
by $\tau$.
The only possibility is that there is a hidden $\tau$
extension from $c_0 d_0$ to $P d_0$.
\end{proof}

\begin{remark}
\label{rem:t-Dh1j1}
If $M \D h_1 d_0$ and $M \D h_1^2 d_0$ are not hit by differentials,
then a straightforward analysis of the sequence
(\ref{eq:LES-Ctau})
shows that the possible $\tau$ extensions on $\D h_1 j_1$ and 
$\tau M h_0 g^2$ must occur.  Thus these uncertainties are entirely
determined by corresponding uncertainties in values of the Adams 
differentials.
\end{remark}

\begin{lemma}
\label{lem:t-d1e1}
\mbox{}
\begin{enumerate}
\item
\revdeg{70, 10, 38}
The element 
$h_1 h_3 (\D e_1 + C_0) + \tau h_2 C''$ maps to
$\ol{h_1^4 c_0 Q_2}$ 
under inclusion of the bottom cell into $C\tau$.
\item
\revdeg{70, 8, 39}
There is a hidden $\tau$ extension from $d_1 e_1$ to 
$h_1 h_3 (\D e_1 + C_0)$.
\end{enumerate}
\end{lemma}

\begin{proof}
Consider the exact sequence
$\pi_{70, 38} \map \pi_{70, 38} C\tau \map \pi_{69, 39}$.
For combinatorial reasons, one of the following two possibilities
must occur:
\begin{enumerate}[(a)]
\item
the element $h_1 h_3 (\D e_1 + C_0) + \tau h_2 C''$ maps to
$\ol{h_1^4 c_0 Q_2}$ 
under inclusion of the bottom cell into $C\tau$, and
there is a hidden $\tau$ extension from $d_1 e_1$ to 
$h_1 h_3 (\D e_1 + C_0)$.
\item
the element $h_1 h_3 (\D e_1 + C_0)$
maps to $\ol{h_1^4 c_0 Q_2}$ 
under inclusion of the bottom cell into $C\tau$, and
there is a hidden $\tau$ extension from $d_1 e_1$ to 
$h_1 h_3 (\D e_1 + C_0) + \tau h_2 C''$.
\end{enumerate}
We will show that there cannot be a hidden
$\tau$ extension from $d_1 e_1$ to
$h_1 h_3 (\D e_1 + C_0) + \tau h_2 C''$.

Lemma \ref{lem:eta-sigma-k1} shows that
$\tau \nu \{d_1 e_1\}$ equals
$\tau \eta \sigma \{k_1\}$.
Since there is no hidden $\tau$ extension on $h_1 k_1$,
there must exist an element $\alpha$ in $\{k_1\}$ such that
$\tau \eta \alpha = 0$.
Therefore,
$\tau \nu \{d_1 e_1\}$ must be zero.

If there were a $\tau$ extension from
$d_1 e_1$ to 
$h_1 h_3 (\D e_1 + C_0) + \tau h_2 C''$, then
$\tau \nu \{d_1 e_1\}$ would be detected by
\[
h_2 \cdot (h_1 h_3 (\D e_1 + C_0) + \tau h_2 C'') =
\tau h_2^2 C'',
\]
and in particular would be non-zero.
\end{proof}

\section{Hidden $2$ extensions}
\label{sctn:2-extn}

\begin{thm}
\label{thm:2-extn}
Tables \ref{tab:2-extn} \rev{and \ref{tab:2-extn-null}} list some hidden extensions by $2$.
\end{thm}

\begin{proof}
Many of the hidden extensions follow by comparison to 
$C\tau$.  For example, there is a hidden $2$ extension
from $h_0 h_2 g$ to $h_1 c_0 d_0$ in the Adams spectral
sequence for $C\tau$.
Pulling back along inclusion of the bottom cell into $C\tau$,
there must also be a hidden $2$ extension from $h_0 h_2 g$
to $h_1 c_0 d_0$ in the Adams spectral sequence for the sphere.
This type of argument is indicated by the notation
$C\tau$ in the fourth column of Table \ref{tab:2-extn}.

Next, Table \ref{tab:tau-extn} shows a hidden $\tau$
extension from $h_1 c_0 d_0$ to $P h_1 d_0$.
Therefore, there is also a hidden $2$ extension from
$\tau h_0 h_2 g$ to $P h_1 d_0$.
This type of argument is indicated by the notation $\tau$
in the fourth column of Table \ref{tab:2-extn}. 

Many cases require more complicated arguments.  In stems up to
approximately dimension 62, see 
\cite{Isaksen14c}*{Section 4.2.2 and Tables 27--28}
\cite{WangXu18}, and \cite{Xu16}.
The higher-dimensional cases are handled in the following lemmas.
\end{proof}

\begin{remark}
Through the 90-stem, there are no crossing $2$ extensions.
\end{remark}

\begin{remark}
The hidden $2$ extension from $h_0 h_3 g_2$ to $\tau g n$
is proved in \cite{WangXu18}, which uses on the
``$\R P^\infty$-method" to establish a hidden
$\sigma$ extension from $\tau h_3 d_1$ to $\D h_2 c_1$
and a hidden $\eta$ extension from $\tau h_1 g_2$ to $\D h_2 c_1$.
We now have easier proofs for these $\eta$ and $\sigma$ extensions,
using the hidden $\tau$ extension from $h_1^2 g_2$ to $\D h_2 c_1$
given in Table \ref{tab:tau-extn}, as well as the relation
$h_3^2 d_1 = h_1^2 g_2$.
\end{remark}

\begin{remark}
\label{rem:2-h0h5i}
Comparison to synthetic homotopy gives additional
information about some
possible hidden $2$ extensions, including:
\begin{enumerate}
\item
there is a hidden $2$ extension from $h_0 h_5 i$ to
$\tau^4 e_0^2 g$.
\item
there is no hidden $2$ extension from $P x_{76,6}$ 
to $M \D h_1 d_0$.
\end{enumerate}
See \rev{\cite{Burklund21} and} \cite{BIX} for more details.
We are grateful to John Rognes for pointing out a mistake
in \cite{Isaksen14c}*{Lemma 4.56 and Table 27} 
concerning the hidden $2$ extension
on $h_0 h_5 i$.  
Lemma \ref{lem:2-h0h5i} shows that the extension occurs but
does not determine its target precisely.  
\end{remark}

\begin{remark}
The first correct proof of the relation $2 \theta_5 = 0$ 
appeared in \cite{Xu16}.  Earlier claims in \cite{Lin01}
and \cite{Kochman90} were based upon a mistaken understanding of the
Toda bracket $\langle \sigma^2, 2, \theta_4 \rangle$.
See \cite{Isaksen14c}*{Table 23} for the correct value of
this bracket.
\end{remark}

\begin{remark}
\label{rem:2-tMh0g^2}
If $M \D h_1^2 d_0$ is non-zero in the $E_\infty$-page, then
there is a hidden $\tau$ extension from $\tau M h_0 g^2$ to
$M \D h_1^2 d_0$.  This implies that there must be a hidden
$2$ extension from $\tau^2 M g^2$ to $M \D h_1^2 d_0$.
\end{remark}

\rev{
\begin{remark}
\label{rem:2-x87,7}
Table \ref{tab:2-extn} shows that there is a hidden $2$
extension from $x_{87,7}$ to $\tau^3 g Q_3$.
This follows from data recently produced by Dexter Chua on the
$d_2$ differentials in the Adams spectral sequence for the cofiber 
of $2$.
\end{remark}
}

\begin{thm}
\label{thm:2-extn-possible}
Table \ref{tab:2-extn-possible} lists all unknown hidden $2$
extensions, through the $90$-stem.
\end{thm}

\begin{proof}
Many possibilities are eliminated by comparison to $C\tau$, to
$\tmf$, or to $\mmf$.  For example,
there cannot be a hidden $2$ extension from $h_2^2 h_4$ to $\tau h_1 g$
by comparison to $C\tau$.

Many additional possibilities are eliminated by consideration of other
parts of the multiplicative structure.  For example,
there cannot be a hidden $2$ extension from $P h_1 h_5$
to $\tau^3 g^2$ because $\tau^3 g^2$ supports an $h_1$ extension
and $2 \eta$ equals zero.

Several cases are a direct consequence of Proposition
\ref{prop:2-<alpha,2,theta5>}.

Some possibilities are eliminated by more complicated arguments.
These cases are handled in the following lemmas.
\end{proof}



\begin{remark}
\label{rem:2-Px76,6}
If $M \D h_1^2 d_0$ is not zero in the $E_\infty$-page, then
$M \D h_1 d_0$ supports an $h_1$ multiplication, and there cannot
be a hidden $2$ extension from $P x_{76,6}$ to $M \D h_1 d_0$.
\end{remark}

\begin{prop}
\label{prop:2-<alpha,2,theta5>}
Suppose that $2 \alpha$ and $\tau \eta \alpha$ are both zero.
Then $2 \langle \alpha, 2, \theta_5 \rangle$ is zero.
\end{prop}

\begin{proof}
Consider the shuffle
\[
2 \langle \alpha, 2, \theta_5 \rangle =
\langle 2, \alpha, 2 \rangle \theta_5.
\]
Since $2 \theta_5$ is zero, this expression has no indeterminacy.
Corollary \ref{cor:2-symmetric} implies that it equals
$\tau \eta \alpha \theta_5$, which is zero by assumption.
\end{proof}

\begin{remark}
\label{rem:2-<alpha,2,theta5>}
Proposition \ref{prop:2-<alpha,2,theta5>} eliminates possible
hidden $2$ extensions on several elements, including $h_2^2 h_6$,
$h_0^3 h_3 h_6$, $h_3^2 h_6$, $h_6 c_1$, $h_2^2 h_4 h_6$,
$h_0^5 h_6 i$, and $h_2^2 h_6 g$.
\end{remark}



\begin{lemma}
\label{lem:2-h0h5i}
\revdeg{54, 9, 28}
There is a hidden $2$ extension from $h_0 h_5 i$ to either
$\tau M P h_1$ or to $\tau^4 e_0^2 g$.
\end{lemma}

\begin{proof}
Table \ref{tab:unit-mmf} shows that
$h_0 h_5 i$ maps to $\D^2 h_2^2$ in the homotopy of $\tmf$.
The element $\D^2 h_2^2$ supports a hidden $2$ extension,
so $h_0 h_5 i$ must support a hidden $2$ extension as well.
\end{proof}

\begin{lemma}
\label{lem:2-th1H1}
\mbox{}
\begin{enumerate}
\item
\revdeg{63, 6, 33}
There is a hidden $2$ extension from
$\tau h_1 H_1$ to $\tau h_1 (\D e_1 + C_0)$.
\item
\revdeg{63, 7, 33}
There is no hidden $2$ extension on $\tau X_2 + \tau C'$.
\item
\revdeg{70, 7, 37}
There is a hidden $2$ extension from
$\tau h_1 h_3 H_1$ to $\tau h_1 h_3 (\D e_1 + C_0)$.
\end{enumerate}
\end{lemma}

\begin{proof}
Table \ref{tab:eta-extn} shows that there is an $\eta$
extension from $\tau h_1 H_1$ to $h_3 Q_2$.
Let $\alpha$ be any element of $\pi_{63,33}$ that is
detected by $\tau h_1 H_1$.
Then
$\tau \eta^2 \alpha$ is non-zero and detected by 
$\tau h_1 h_3 Q_2$.
Note that $\tau h_1 h_3 Q_2$ cannot be the target of a hidden
$2$ extension because there are no possibilities.

If $2 \alpha$ were zero, then we would have the shuffling
relation
\[
\tau \eta^2 \alpha =
\langle 2, \eta, 2 \rangle \alpha = 
2 \langle \eta, 2, \alpha \rangle.
\]
But this would contradict the previous paragraph.

We now know that $2 \alpha$ must be non-zero for every possible
choice of $\alpha$.
The only possibility is that there is a hidden
$2$ extension from $\tau h_1 H_1$ to $\tau h_1 (\D e_1 + C_0)$,
and that there is no hidden $2$ extension on $\tau X_2 + \tau C'$.
This establishes the first two parts.

The third part follows immediately from the first part by
multiplication by $h_3$.
\end{proof}

\begin{lemma}
\label{lem:2-h1h6}
\revdeg{64, 2, 33}
There is a hidden $2$ extension from $h_1 h_6$ to $\tau h_1^2 h_5^2$.
\end{lemma}

\begin{proof}
Table \ref{tab:Toda} shows that 
$h_1 h_6$ detects $\langle \eta, 2, \theta_5 \rangle$.
Now shuffle to obtain
\[
2 \langle \eta, 2, \theta_5 \rangle =
\langle 2, \eta, 2 \rangle \theta_5 = \tau \eta^2 \theta_5.
\]
\end{proof}

\begin{lemma}
\label{lem:2-D1h3^2}
\revdeg{66, 6, 36}
There is no hidden $2$ extension on $\D_1 h_3^2$.
\end{lemma}

\begin{proof}
Table \ref{tab:Toda} shows that $\D_1 h_3^2$ detects the
Toda bracket $\langle \eta^2, \theta_4, \eta^2, \theta_4 \rangle$.
We have
\[
2 \langle \eta^2, \theta_4, \eta^2, \theta_4 \rangle \subseteq
\langle \langle 2, \eta^2, \theta_4 \rangle, \eta^2, \theta_4 \rangle.
\]
Table \ref{tab:Toda} shows that
\[
\nu \theta_4 = \langle 2, \eta, \eta \theta_4 \rangle =
\langle 2, \eta^2, \theta_4 \rangle,
\]
so we must compute $\langle \nu \theta_4, \eta^2, \theta_4 \rangle$.

This bracket contains $\nu \langle \theta_4, \eta^2, \theta_4 \rangle$,
which equals zero by Lemma \ref{lem:theta4,eta^2,theta4}.
Therefore, we only need to compute the indeterminacy of 
$\langle \nu \theta_4, \eta^2, \theta_4 \rangle$.

The only possible non-zero element in the indeterminacy is the 
product $\theta_4 \{ t \}$.
Table \ref{tab:Toda} shows that 
$\{t \} = \langle \nu, \eta, \eta \theta_4 \rangle$.  Now
\[
\theta_4 \{ t\} = \langle \nu, \eta, \eta \theta_4 \rangle \theta_4 =
\nu \langle \eta, \eta \theta_4, \theta_4 \rangle.
\]
This last expression is well-defined because
$\theta_4^2$ is zero \cite{Xu16}, and it 
must be zero because 
$\pi_{63,34}$ consists entirely of multiples of $\eta$.
\end{proof}

\begin{lemma}
\label{lem:2-h0Q3}
\revdeg{67, 6, 36}
There is no hidden $2$ extension on $h_0 Q_3 + h_2^2 D_3$.
\end{lemma}

\begin{proof}
By comparison to the homotopy of $C\tau$,
there is no hidden extension with value $h_2^2 A'$.
Table \ref{tab:eta-extn} shows that
$\tau^2 \D h_2^2 e_0 g$ supports a hidden $\eta$ extension.
Therefore, it cannot be the target of a $2$ extension.
\end{proof}

\begin{lemma}
\label{lem:2-h3A'}
\revdeg{68, 7, 36}
There is no hidden $2$ extension on $h_3 A'$.
\end{lemma}

\begin{proof}
Table \ref{tab:Toda} shows that $h_3 A'$ detects the Toda bracket
$\langle \sigma, \kappa, \tau \eta \theta_{4.5} \rangle$.
Shuffle to obtain
\[
\langle \sigma, \kappa, \tau \eta \theta_{4.5} \rangle 2 =
\sigma \langle \kappa, \tau \eta \theta_{4.5}, 2 \rangle.
\]
The bracket $\langle \kappa, \tau \eta \theta_{4.5}, 2 \rangle$
is zero because it is contained in $\pi_{61,32} = 0$.
\end{proof}


\begin{lemma}
\label{lem:2-p'}
\revdeg{69, 4, 36}
There is no hidden $2$ extension on $p'$.
\end{lemma}

\begin{proof}
Table \ref{tab:misc-extn} shows that $p'$
detects the product $\sigma \theta_5$, and 
$2 \theta_5$ is already known to be zero
\cite{Xu16}.
\end{proof}

\begin{lemma}
\label{lem:2-th1D3'}
\revdeg{70, 9, 37}
There is no hidden $2$ extension on $\tau h_1 D_3'$.
\end{lemma}

\begin{proof}
Table \ref{tab:Toda} shows that 
$\tau h_1 D'_3$ detects the Toda bracket 
$\langle \eta, \nu, \tau \theta_{4.5} \kappabar \rangle$.
Now shuffle to obtain
\[
2 \langle \eta, \nu, \tau \theta_{4.5} \kappabar \rangle =
\langle 2, \eta, \nu \rangle \tau \theta_{4.5} \kappabar,
\]
which equals zero because $\langle 2, \eta, \nu \rangle$
is contained in in $\pi_{5,3} = 0$.
\end{proof}

\begin{lemma}
\label{lem:2-h1h3h6}
\revdeg{71, 3, 37}
There is no hidden $2$ extension on $h_1 h_3 h_6$.
\end{lemma}

\begin{proof}
Table \ref{tab:Toda} shows that
$h_1 h_6$ detects the Toda bracket $\langle \eta, 2, \theta_5 \rangle$.
Let $\alpha$ be an element of this bracket.
Then $h_1 h_3 h_6$ detects $\sigma \alpha$, and
\[
2 \sigma \alpha = 2 \sigma \langle \eta, 2, \theta_5 \rangle =
\sigma \langle 2, \eta, 2 \rangle \theta_5 = 
\tau \eta^2 \sigma \theta_5.
\]

Table \ref{tab:Toda} also shows that $h_6 c_0$ detects 
the Toda bracket $\langle \epsilon, 2, \theta_5 \rangle$.
Let $\beta$ be an element of this bracket.
As in the proof of Lemma \ref{lem:2-h6c0}, we compute
that $2 \beta$ equals $\tau \eta \epsilon \theta_5$.

Now consider the element $\sigma \alpha + \beta$, which is also
detected by $h_1 h_3 h_6$.  Then
\[
2(\sigma \alpha + \beta) = 
\tau \eta^2 \sigma \theta_5 + \tau \eta \epsilon \theta_5 =
\tau \nu^3 \theta_5,
\]
using Toda's relation $\eta^2 \sigma + \nu^3 = \eta \epsilon$
\cite{Toda62}.

Table \ref{tab:nu-extn} shows that there is a hidden
$\nu$ extension from $h_2 h_5^2$ to $\tau h_1 Q_3$.
Therefore, $\tau h_1 Q_3$ detects $\nu^2 \theta_5$.

This does not yet imply that $\nu^3 \theta_5$ is zero,
because $\nu^2 \theta_5 + \eta \{\tau Q_3 + \tau n_1\}$
might be detected $h_3 A'$ or $P h_2 h_5 j$ in higher filtration.
However, $h_3 A'$ does not support a hidden
$\nu$ extension by Lemma \ref{lem:nu-h3A'}.
Also, Table \ref{tab:unit-mmf} shows that
$P h_2 h_5 j$ maps non-trivially to $\tmf$,
while $\nu^2 \theta_5 + \eta \{\tau Q_3 + \tau n_1\}$
maps to zero.
This is enough to conclude that $\nu^3 \theta_5$ is zero.

We have now shown that $2(\sigma \alpha + \beta)$ is zero
in $\pi_{71,37}$.  Since $h_1 h_3 h_6$
detects $\sigma \alpha + \beta$, it follows that 
$h_1 h_3 h_6$ does not support a hidden $2$ extension.
\end{proof}

\begin{lemma}
\label{lem:2-h6c0}
\revdeg{71, 4, 37}
There is a hidden $2$ extension from $h_6 c_0$ to $\tau h_1^2 p'$.
\end{lemma}

\begin{proof}
Table \ref{tab:Toda} shows that $h_6 c_0$ detects the Toda bracket
$\langle \epsilon, 2, \theta_5 \rangle$.
Now shuffle to obtain
\[
2 \langle \epsilon, 2, \theta_5 \rangle = 
\langle 2, \epsilon, 2 \rangle \theta_5 = 
\tau \eta \epsilon \theta_5.
\]
Finally, $\tau \eta \epsilon \theta_5$ is detected
by $\tau h_1^2 p'$ because of the relation
$h_1^2 p' = h_1 h_5^2 c_0$.
\end{proof}

\begin{lemma}
\label{lem:2-th1p1}
\revdeg{71, 5, 37}
There is no hidden $2$ extension on $\tau h_1 p_1$.
\end{lemma}

\begin{proof}
Lemma \ref{lem:eta,nu,t^2h2C'} shows that $\tau h_1 p_1$
detects $\langle \eta, \nu, \alpha \rangle$ for some
$\alpha$ detected by $\tau^2 h_2 C'$.
Now shuffle to obtain
\[
2 \langle \eta, \nu, \alpha \rangle = 
\langle 2, \eta, \nu \rangle \alpha,
\]
which is zero because $\langle 2, \eta, \nu \rangle$ is contained
in $\pi_{5,3} = 0$.
\end{proof}

\begin{lemma}
\label{lem:2-h1^2h3H1}
\revdeg{71, 8, 39}
There is a hidden $2$ extension from $h_2^3 H_1$ to $\tau M h_2^2 g$.
\end{lemma}

\begin{proof}
Table \ref{tab:nu-extn} shows that there are hidden $\nu$ extensions
from $h_2^3 H_1$ to $h_3 C''$, and from
$\tau M h_2^2 g$ to $M h_1 d_0^2$.
Table \ref{tab:2-extn} shows that there is also a hidden
$2$ extension from $h_3 C''$ to $M h_1 d_0^2$.
The only possibility is that there must also be a hidden $2$
extension on $h_1^2 h_3 H_1$.
\end{proof}

\begin{lemma}
\label{lem:2-Ph1h6}
\revdeg{72, 6, 37}
There is no hidden $2$ extension on $P h_1 h_6$.
\end{lemma}

\begin{proof}
Table \ref{tab:Toda} shows that $P h_1 h_6$ detects the Toda bracket
$\langle \mu_9, 2, \theta_5 \rangle$.
Shuffle to obtain
\[
2 \langle \mu_9, 2, \theta_5 \rangle =
\langle 2, \mu_9, 2 \rangle \theta_5 =
\tau \eta \mu_9 \theta_5.
\]
Table \ref{tab:Toda} also shows that $\mu_9$ is contained in the
Toda bracket $\langle \eta, 2, 8 \sigma \rangle$.
Shuffle again to obtain
\[
\tau \eta \mu_9 \theta_5 =
\langle \eta, 2, 8 \sigma \rangle \tau \eta \theta_5 =
\tau \eta^2 \langle 2, 8 \sigma, \theta_5 \rangle.
\]
Table \ref{tab:Toda} shows that
$h_0^3 h_3 h_6$ detects $\langle 2, 8 \sigma, \theta_5 \rangle$.

By inspection, the product $\eta^2 \{h_0^3 h_3 h_6\}$
can only be detected by $\D^2 h_1 h_4 c_0$.
However, this cannot occur by comparison to $C\tau$.
Therefore,
$\eta^2 \{h_0^3 h_3 h_6\}$, and also
$\tau \eta^2 \{h_0^3 h_3 h_6\}$, must be zero.
\end{proof}

\begin{lemma}
\label{lem:2-h4Q2+h3^2D2}
\revdeg{72, 8, 38}
There is no hidden $2$ extension on $h_4 Q_2 + h_3^2 D_2$.
\end{lemma}

\begin{proof}
Table \ref{tab:Toda} shows that the element
$h_4 Q_2 + h_3^2 D_2$ detects the Toda bracket
$\langle \sigma^2, 2, \{t\}, \tau \kappabar \rangle$.
Consider the relation
\[
2 \langle \sigma^2, 2, \{t\}, \tau \kappabar \rangle \subseteq
\langle \langle 2, \sigma^2, 2 \rangle, \{t\}, \tau \kappabar \rangle.
\]
Corollary \ref{cor:2-symmetric} shows 
that the Toda bracket $\langle 2, \sigma^2, 2 \rangle$ contains
zero since $\eta \sigma^2$ is zero.  Therefore, it
consists of even multiples of $\rho_{15}$;
let $2 k \rho_{15}$ be any such element in $\pi_{15,8}$.

The Toda bracket
$\langle 2 k \rho_{15}, \{t\}, \tau \kappabar \rangle$
contains $k \rho_{15} \langle 2, \{t\}, \tau \kappabar \rangle$,
which equals zero as discussed in the proof
of Lemma \ref{lem:sigma^2,2,T,tkappabar}.
Moreover, its indeterminacy
is equal to $\tau \kappabar \cdot \pi_{52,28}$,
which is detected in Adams filtration at least 12.
This implies that
$\langle 2 k \rho_{15}, \{t\}, \tau \kappabar \rangle$
is detected in Adams filtration at least 12,
and that
the target of a hidden $2$ extension on $h_4 Q_2 + h_3^2 D_2$
must have Adams filtration at least 12.

The remaining possible targets with Adams filtration at least 12
are eliminated by comparison to $C\tau$ or to $\mmf$.
\end{proof}

\begin{remark}
\label{rem:2-h4Q2+h3^2D2}
The proof of Lemma \ref{lem:2-h4Q2+h3^2D2} might be simplified by
considering the shuffle
\[
2 \langle \sigma^2, 2, \{t\}, \tau \kappabar \rangle =
\langle 2, \sigma^2, 2, \{t\} \rangle \tau \kappabar.
\]
However, the latter four-fold bracket may not exist, since
both three-fold subbrackets have indeterminacy.
See \cite{Isaksen14b} for a discussion of the analogous 
difficulty with Massey products.
\end{remark}

\begin{lemma}
\label{lem:2-h2^2Q3}
\revdeg{73, 7, 40}
There is no hidden $2$ extension on $h_2^2 Q_3$.
\end{lemma}

\begin{proof}
The element $\tau h_2^2 Q_3$ detects
$\nu^2 \{\tau Q_3 + \tau n_1\}$, so it cannot support a 
hidden $2$ extension.
This rules out all possible $2$ extensions on
$h_2^2 Q_3$.
\end{proof}

\begin{lemma}
\label{lem:2-h0h4D2}
\revdeg{73, 8, 38}
There is no hidden $2$ extension on $h_0 h_4 D_2$.
\end{lemma}

\begin{proof}
Table \ref{tab:tau-extn} shows that there is a hidden
$\tau$ extension from $h_1^2 h_6 c_0$ to $h_0 h_4 D_2$.
Therefore, $h_0 h_4 D_2$ detects either
$\tau \eta \epsilon \eta_6$ or
$\tau \eta \epsilon \eta_6 + \nu^2 \{\tau Q_3 + \tau n_1\}$,
because of the presence of $\tau h_2^2 Q_3$ in higher
filtration.
In either case, $h_0 h_4 D_2$ cannot support a hidden $2$ extension.
\end{proof}

\begin{lemma}
\label{lem:2-h3(tQ3+tn1)}
\revdeg{74, 6, 39}
There is a hidden $2$ extension from $h_3(\tau Q_3+ \tau n_1)$
to $\tau x_{74,8}$.
\end{lemma}

\begin{proof}
Table \ref{tab:misc-extn} shows that
$\tau x_{74,8}$ detects $\tau \kappabar_2 \theta_4$.
Table \ref{tab:Toda} shows that $\theta_4$ equals
the Toda bracket $\langle \sigma^2, 2, \sigma^2, 2 \rangle$.

Now consider the shuffle
\[
\tau \kappabar_2 \theta_4 = 
\tau \kappabar_2 \langle \sigma^2, 2, \sigma^2, 2 \rangle =
\langle \tau \kappabar_2, \sigma^2, 2, \sigma^2 \rangle 2.
\]
Lemma \ref{lem:tkappabar2,sigma^2,2} shows that the latter
bracket is well-defined.
This implies that $\tau x_{74,8}$ is the target of a hidden $2$
extension, and $h_3(\tau Q_1 + \tau n_1)$ is the only possible
source.
\end{proof}


\begin{lemma}
\label{lem:2-h6d0}
\revdeg{77, 5, 40}
There is no hidden $2$ extension on $h_6 d_0$.
\end{lemma}

\begin{proof}
Table \ref{tab:Toda} shows that $h_6 d_0$ detects the Toda bracket
$\langle \kappa, 2, \theta_5 \rangle$.  Now shuffle to obtain
\[
2 \langle \kappa, 2, \theta_5 \rangle = 
\langle 2, \kappa, 2 \rangle \theta_5 = \tau \eta \kappa \theta_5.
\]
Lemma \ref{lem:sigma^2+kappa-h5^2} shows that this product
equals \rev{either $\tau \eta \sigma^2 \theta_5$ or
$\tau \eta \sigma^2 \theta_5 + \tau^3 \eta \kappa_1 \kappabar_2$.
Both possibilities are zero because
$\eta \sigma^2$ and $\tau \eta \kappa_1$ are zero.}
\end{proof}

\begin{lemma}
\label{lem:2-e0A'}
\revdeg{78, 10, 42}
There is a hidden $2$ extension from $e_0 A'$ to $M \D h_1^2 h_3$.
\end{lemma}

\begin{proof}
Let $\alpha$ be an element of $\pi_{76,40}$ that is detected
by $x_{76,9}$.
Table \ref{tab:eta-extn} shows that there is a hidden
$\eta$ extension from $x_{76,9}$ to $M \D h_1 h_3$,
so $\tau \eta^2 \alpha$ is detected by $\tau M \D h_1^2 h_3$.
Now shuffle to obtain
\[
\tau \eta^2 \alpha = \langle 2, \eta, 2 \rangle \alpha =
2 \langle \eta, 2, \alpha \rangle.
\]
This shows that $\tau M \D h_1^2 h_3$ must be the target
of a hidden $2$ extension.

Moreover, the source of this hidden $2$ extension must be in Adams
filtration at least $10$, since the Adams differential
$d_2(\tau x_{77,8}) = h_0 x_{76,9}$ implies that
$\langle \eta, 2, \alpha \rangle$ is detected by
$h_1 x_{77,8} = 0$ in filtration $9$.
The only possible source is $e_0 A'$.
\end{proof}

\begin{lemma}
\label{lem:2-h1h4h6}
\revdeg{79, 3, 41}
There is no hidden $2$ extension on $h_1 h_4 h_6$.
\end{lemma}

\begin{proof}
Table \ref{tab:Toda} shows that $h_1 h_4 h_6$ detects the Toda bracket
$\langle \eta_4, 2, \theta_5 \rangle$.
Now shuffle to obtain
\[
2 \langle \eta_4, 2, \theta_5 \rangle =
\langle 2, \eta_4, 2 \rangle \theta_5,
\]
which equals $\tau \eta \eta_4 \theta_5$ by Table \ref{tab:Toda}.
We will show that this product is zero.

There are several elements in the Adams $E_\infty$-page
that might detect $\eta_4 \theta_5$.  The possibilities
$h_1 h_6 d_0$ and $x_{78,9}$ are ruled out by comparison to $C\tau$.
The possibility $\tau e_0 A'$ is ruled out because
Table \ref{tab:2-extn} shows that $e_0 A'$ supports a hidden
$2$ extension.

Two possibilities remain.  If $\eta_4 \theta_5$ is detected
by $\tau M \D h_1^2 h_3$, then $\tau \eta \eta_4 \theta_5$
must be zero because there are no elements in sufficiently
high Adams filtration.

Finally, suppose that $\eta_4 \theta_5$ is detected by
$\tau h_1^2 x_{76,6}$.  Let $\alpha$ be an element of $\pi_{77,41}$
that is detected by $h_1 x_{76,6}$.
If $\eta_4 \theta_5 + \tau \eta \alpha$ is not zero,
then it is detected in
higher filtration.  
\rev{It cannot be detected by $x_{78,9}$ by comparison to $C\tau$, and}
it cannot be detected by $\tau e_0 A'$
because of the hidden $2$ extension on $e_0 A'$.
If it is detected by $\tau M \D h_1^2 h_3$, then we may
change the choice of $\alpha$ to ensure that
$\eta_4 \theta_5 + \tau \eta \alpha$ is zero.

We have now shown that $\tau \eta \eta_4 \theta_5$
equals $\tau^2 \eta^2 \alpha$.
Shuffle to obtain
\[
\tau^2 \eta^2 \alpha =
\tau \alpha \langle 2, \eta, 2 \rangle =
\langle \alpha, 2, \eta \rangle 2 \tau.
\]
\rev{Here we are using that $2 \alpha$ is zero; the possible
$2$ extensions on $h_1 x_{76,6}$ are easily eliminated by the presence
of $\eta$ extensions and by comparison to $\mmf$.}

Table \ref{tab:Toda} shows that
$\langle \alpha, 2, \eta \rangle$ is detected
by $h_0 h_2 x_{76,6}$, and Lemma \ref{lem:2-h0h2x76,6}
shows that this element does not support a hidden $2$ extension.
Therefore,
$\langle \alpha, 2, \eta \rangle 2 \tau$ is zero.
\end{proof}

\begin{lemma}
\label{lem:2-h0h2x76,6}
\revdeg{79, 8, 42}
There is no hidden $2$ extension on $h_0 h_2 x_{76,6}$.
\end{lemma}

\begin{proof}
Let $\alpha$ be an element of $\pi_{76,40}$ that is detected by
$h_4 A$.  Then $\nu \alpha$ is detected by $h_0 h_2 x_{76,6}$,
\rev{and $2 \alpha$ is detected by $h_0 h_4 A$.
We will show that $2 \nu \alpha$ is zero.}

Table \ref{tab:misc-extn} shows that
$h_0 h_4 A$ also detects 
\rev{either $\sigma^2 \theta_5$ or $\sigma^2 \theta_5 + \tau^2 \kappa_1 \kappabar_2$.
Then
$2 \alpha + \sigma^2 \theta_5$ or $2 \alpha + \sigma^2 \theta_5 + \tau^2 \kappa_1 \kappabar_2$ could be detected in higher filtration.
However, only $x_{76,9}$ could detect this error term, and
inclusion of the bottom cell into $C\tau$ rules it out.

We now know that
$\sigma^2 \theta_5 + 2 \alpha$ is either zero or
$\tau^2 \kappa_1 \kappabar_2$.
Multiply by $\nu$ to conclude that
$2 \nu \alpha$ is either zero or $\tau^2 \nu \kappa_1 \kappabar_2$.
As in the proof of Lemma \ref{lem:eta-h1h6d0}, this last
expression is also zero.
}
\end{proof}

\begin{lemma}
\label{lem:2-Ph6c0}
\revdeg{79, 8, 41}
There is no hidden $2$ extension on $P h_6 c_0$.
\end{lemma}

\begin{proof}
Table \ref{tab:misc-extn} shows that
$P h_6 c_0$ detects the product $\rho_{15} \eta_6$.
Table \ref{tab:Toda} shows that 
$\eta_6$ is contained in the Toda bracket
$\langle \eta, 2, \theta_5 \rangle$.
Now shuffle to obtain
\[
2 \rho_{15} \eta_6 =
2 \rho_{15} \langle \eta, 2, \theta_5 \rangle =
\rho_{15} \theta_5 \langle 2, \eta, 2 \rangle,
\]
which equals $\tau \eta^2 \rho_{15} \theta_5$ by 
Table \ref{tab:Toda}.

Table \ref{tab:misc-extn} shows that
$\rho_{15} \theta_5$ is detected by either
$h_0 x_{77,7}$ or $\tau^2 m_1$.
First suppose that it is detected by $h_0 x_{77,7}$.
Table \ref{tab:2-extn} shows that $h_0 x_{77,7}$ is the target
of a $2$ extension.
Then $\rho_{15} \theta_5$ equals $2 \alpha$ modulo
higher filtration.
In any case,
$\tau \eta^2 \rho_{15} \theta_5$ is zero.

Next suppose that $\rho_{15} \theta_5$ is detected by
$\tau^2 m_1$.
Then $\rho_{15} \theta_5$ equals $\tau^2 \alpha$ modulo
higher filtration for some element $\alpha$ detected by $m_1$.
Table \ref{tab:tau-extn}
shows that there is a hidden $\tau$ extension from $h_1 m_1$
to $M \D h_1^2 h_3$.  This implies that
$\tau \eta \alpha$ is detected by $M \D h_1^2 h_3$.
Finally,
$\tau^3 \eta^2 \alpha = \tau \eta^2 \rho_{15} \theta_5$
must be zero.
\end{proof}

\begin{lemma}
\label{lem:2-DB6}
\revdeg{79, 11, 42}
There is no hidden $2$ extension on $\D B_6$.
\end{lemma}

\begin{proof}
Table \ref{tab:eta-extn} shows that there is a hidden $\eta$ extension
from $h_0^6 h_4 h_6$ to $\tau \D B_6$.
Therefore, $\tau \D B_6$ cannot be the source of a hidden $2$
extension, so there cannot be a hidden $2$ extension
from $\D B_6$ to $\tau^2 M e_0^2$.
\end{proof}

\rev{
\begin{lemma}
\label{lem:2-h5^2g}
\revdeg{82, 6, 44}
There is no hidden $2$ extension on $h_5^2 g$.
\end{lemma}

\begin{proof}
The element $\tau h_5^2 g$ detects the product
$\kappabar \theta_5$, so it cannot support a hidden $2$
extension since $2 \theta_5$ is zero.

If there were a hidden $2$ extension from
$h_5^2 g$ to $\tau (\D e_1 + C_0) g$, then
the hidden 
$\tau$ extension from $\tau (\D e_1 + C_0) g$ to $\D^2 h_2 n$
would imply that
there is a hidden $2$ extension from 
$\tau h_5^2 g$ to $\D^2 h_2 n$.
\end{proof}
}

\begin{lemma}
\label{lem:2-h2^2x76,6}
\revdeg{82, 8, 44}
There is no hidden $2$ extension on $h_2^2 x_{76,6}$.
\end{lemma}

\begin{proof}
As in the proof of Lemma \ref{lem:2,eta,h2x76,6},
let $\alpha$ be an element in $\pi_{79,42}$ that is detected by
$h_2 x_{76,6}$ such that $\eta \alpha$ is zero.
Then $\nu \alpha$ is detected by $h_2^2 x_{76,6}$, and we wish to show
that $2 \nu \alpha$ is zero.

Table \ref{tab:Toda} shows that $2 \nu$ is contained in
$\langle \eta, 2, \eta \rangle$.
Consider the shuffle
\[
2 \nu \alpha = \langle \eta, 2, \eta \rangle \alpha =
\eta \langle 2, \eta, \alpha \rangle.
\]
Table \ref{tab:Toda} shows that the last Toda bracket is zero.
\end{proof}

\begin{lemma}
\label{lem:2-h0^2h6g}
\revdeg{83, 7, 44}
There is no hidden $2$ extension on $h_0^2 h_6 g$.
\end{lemma}

\begin{proof}
The element $h_0^2 h_6 g$ equals
$h_2^2 h_6 d_0$, so it detects $\nu^2 \{h_6 d_0\}$.
\end{proof}

\begin{lemma}
\label{lem:2-th2h4Q3}
\revdeg{85, 7, 45}
There is no hidden $2$ extension on $\tau h_2 h_4 Q_3$.
\end{lemma}

\begin{proof}
There cannot be a hidden $2$ extension from $\tau h_2 h_4 Q_3$
to $\tau P h_1 x_{76,6}$ because there is no hidden
$\tau$ extension from $h_0 h_2 h_4 Q_3$ to $P h_1 x_{76,6}$.

Table \ref{tab:eta-extn} shows that 
$\tau^3 M g^2$ supports a hidden $\eta$ extension.  Therefore, it
cannot be the target of a hidden $2$ extension.
\end{proof}

\begin{lemma}
\label{lem:2-h6hc0d0}
\mbox{}
\begin{enumerate}
\item
\revdeg{85, 8, 45}
There is no hidden $2$ extension on $h_6 c_0 d_0$.
\item
\revdeg{85, 9, 44}
There is no hidden $2$ extension on $P h_6 d_0$.
\end{enumerate}
\end{lemma}

\begin{proof}
Table \ref{tab:eta-extn} shows that both elements
are targets of hidden $\eta$ extensions.
\end{proof}

\begin{lemma}
\label{lem:2-h4h6c0}
\revdeg{86, 5, 45}
There is no hidden $2$ extension on $h_4 h_6 c_0$.
\end{lemma}

\begin{proof}
Table \ref{tab:misc-extn} shows that
$h_4 h_6 c_0$ detects the product $\sigma \{h_1 h_4 h_6\}$,
and the element $h_1 h_4 h_6$ does not support a hidden $2$ extension
by Lemma \ref{lem:2-h1h4h6}.
\end{proof}

\begin{lemma}
\label{lem:2-th2gC'}
\revdeg{86, 12, 47}
There is no hidden $2$ extension on $\tau h_2 g C'$.
\end{lemma}

\begin{proof}
The possible target $\tau^3 e_0^3 m$ is ruled out by comparison
to $\mmf$.  The possible target $P h_1^7 h_6$ is ruled out by 
comparison to $C\tau$.

It remains to eliminate the possible target $\tau^2 M h_1 g^2$.
Table \ref{tab:tau-extn} shows that there are hidden $\tau$
extensions from
$\tau h_2 g C'$ and $\tau^2 M h_1 g^2$ to
$\D^2 h_2^2 d_1$ and $M \D h_0^2 e_0$ respectively.
However, there is no hidden $2$ extension from
$\D^2 h_2^2 d_1$ to $M \D h_0^2 e_0$, so there cannot be a
$2$ extension from $\tau h_2 g C'$ to $\tau^2 M h_1 g^2$.
\end{proof}

\begin{lemma}
\label{lem:2-h1^2c3}
\revdeg{87, 5, 46}
There is no hidden $2$ extension on $h_1^2 c_3$.
\end{lemma}

\begin{proof}
Table \ref{tab:Toda} shows that the Toda
bracket $\langle \tau \{h_0 Q_3 + h_0 n_1\}, \nu_4, \eta \rangle$
is detected by $h_1^2 c_3$.
Shuffle to obtain
\[
\langle \tau \{h_0 Q_3 + h_0 n_1\}, \nu_4, \eta \rangle 2 =
\tau \{h_0 Q_3 + h_0 n_1\} \langle \nu_4, \eta, 2 \rangle.
\]
These expressions have no indeterminacy because
$\tau \{h_0 Q_3 + h_0 n_1\}$ does not support a $2$ extension.
Finally, the bracket $\langle \nu_4, \eta, 2 \rangle$
contains zero by comparison to $\tmf$.
\end{proof}

\begin{lemma}
\label{lem:2-P^2h6c0}
\revdeg{87, 12, 45}
There is no hidden $2$ extension on $P^2 h_6 c_0$.
\end{lemma}

\begin{proof}
Table \ref{tab:misc-extn} shows that $P^2 h_6 c_0$ detects
the product $\rho_{23} \eta_6$.
Table \ref{tab:Toda} shows that $\eta_6$ is contained in the
Toda bracket $\langle \eta, 2, \theta_5 \rangle$.
Shuffle to obtain
\[
2 \rho_{23} \eta_6 =
2 \rho_{23} \langle \eta, 2, \theta_5 \rangle =
\langle 2, \eta, 2 \rangle \rho_{23} \theta_5 =
\tau \eta^2 \rho_{23} \theta_5.
\]
The product $\rho_{23} \theta_5$ is detected in Adams
filtration at least 13, and then $\tau \eta^2 \rho_{23} \theta_5$
is detected in filtration at least 16.  This rules out all possible
targets for a hidden $2$ extension on $P^2 h_6 c_0$.
\end{proof}


\begin{lemma}
\label{lem:2-M^2}
\revdeg{90, 12, 48}
There is no hidden $2$ extension on $M^2$.
\end{lemma}

\begin{proof}
Table \ref{tab:misc-extn} shows that 
$M^2$ detects $\theta_{4.5}^2$.
Graded commutativity implies that $2 \theta_{4.5}^2$ is zero.
\end{proof}

\section{Hidden $\eta$ extensions}
\label{sctn:eta-extn}

\begin{thm}
\label{thm:eta-extn}
Tables \ref{tab:eta-extn} \rev{and \ref{tab:eta-extn-null}} list some hidden extensions by $\eta$.
\end{thm}

\begin{proof}
Many of the hidden extensions follow by comparison to 
$C\tau$.  For example, there is a hidden $\eta$ extension
from $\tau h_1 g$ to $c_0 d_0$ in the Adams spectral
sequence for $C\tau$.
Pulling back along inclusion of the bottom cell into $C\tau$,
there must also be a hidden $\eta$ extension from $\tau h_1 g$
to $c_0 d_0$ in the Adams spectral sequence for the sphere.
This type of argument is indicated by the notation
$C\tau$ in the fourth column of Table \ref{tab:eta-extn}.

Next, Table \ref{tab:tau-extn} shows a hidden $\tau$
extension from $c_0 d_0$ to $P d_0$.
Therefore, there is also a hidden $\eta$ extension from
$\tau^2 h_1 g$ to $P d_0$.
This type of argument is indicated by the notation $\tau$
in the fourth column of Table \ref{tab:eta-extn}.

The proofs of several of the extensions in Table \ref{tab:eta-extn}
rely on analogous extensions in $\mmf$.  
Extensions in $\mmf$ have not been rigorously analyzed \cite{Isaksen18}.
However, the specific extensions from $\mmf$ that we need
are easily deduced from extensions in $\tmf$, together
with the multiplicative structure.
For example, there is a hidden $\eta$ extension in $\tmf$
from $a n$ to $\tau d_0^2$.
Therefore, there is a hidden $\eta$ extension in $\mmf$
from $a n g$ to $\tau d_0^2 g$, and also a hidden
$\eta$ extension from $\D h_2^2 e_0$ to $\tau d_0 e_0^2$
in the homotopy groups of the sphere spectrum.
Note that $\mmf$ really is required here, since 
$a n g$ and $d_0^2 g$ equal
zero in the homotopy of $\tmf$.


Many cases require more complicated arguments.  In stems up to
approximately dimension 62, see 
\cite{Isaksen14c}*{Section 4.2.3 and Tables 29--30} and
\cite{WangXu18}.
The higher-dimensional cases are handled in the following lemmas.
\end{proof}

\begin{remark}
\label{rem:eta-tC}
The hidden $\eta$ extension from $\tau C$ to $\tau^2 g n$
is proved in \cite{WangXu18}, which uses on the
``$\R P^\infty$-method" to establish a hidden
$\sigma$ extension from $\tau h_3 d_1$ to $\D h_2 c_1$
and a hidden $\eta$ extension from $\tau h_1 g_2$ to $\D h_2 c_1$.
We now have easier proofs for these $\eta$ and $\sigma$ extensions,
using the hidden $\tau$ extension from $h_1^2 g_2$ to $\D h_2 c_1$
given in Table \ref{tab:tau-extn}, as well as the relation
$h_3^2 d_1 = h_1^2 g_2$.
\end{remark}



\rev{
\begin{remark}
\label{rem:eta-tDj1+t^2gC'}
If $h_1 f_2$ survives, then there is a hidden
$\tau$ extension from $\D h_1 j_1$ to $M \D h_1 d_0$.
It follows that there must be a hidden
$\eta$ extension from $\tau \D j_1 + \tau^2 g C'$ to
$M \D h_1 d_0$.
\end{remark}
}

\begin{remark}
The last column of Table \ref{tab:eta-extn}
indicates the crossing $\eta$ extensions.
\end{remark}

\begin{thm}
\label{thm:eta-extn-possible}
Table \ref{tab:eta-extn-possible}
lists all unknown hidden $\eta$ extensions, through
the $90$-stem. 
\end{thm}

\begin{proof}
Many possible extensions can be eliminated by
comparison to $C\tau$, to $\tmf$, or to $\mmf$.
For example,
there cannot be a hidden $\eta$ extension
from $\tau M d_0$ to $\tau^4 g^3$ because
$\tau^4 g^3$ maps to a non-zero element
in $\pi_{60} \tmf$ that is not divisible by $\eta$.

Other possibilities are eliminated by consideration of other
parts of the multiplicative structure.  For example,
there cannot be a hidden $\eta$ extension whose target supports
a multiplication by $2$, since $2 \eta$ equals zero.

Many cases are eliminated by more complicated arguments.
These are handled in the following lemmas.
\end{proof}

\rev{
\begin{remark}
\label{rem:eta-h2D3}
There is no hidden $\eta$ extension on $h_2 D_3$.
The possible target $\tau k_1$ is eliminated by computer
data recently produced by Dexter Chua on $d_2$ differentials
the Adams spectral sequence for the cofiber of $\eta$.
\end{remark}
}

\begin{remark}
\label{rem:eta-tgD3}
There is a hidden $\tau$ extension from
$\tau (\D e_1 + C_0) g$ to $\D^2 h_2 n$.
The possible extension from $\tau g D_3$ to 
$\tau (\D e_1 + C_0) g$ occurs if and only if
the possible extension from $\tau^2 g D_3$ to
$\D^2 h_2 n$ occurs.
\end{remark}

\rev{
\begin{remark}
\label{rem:eta-to-gQ3}
Computer
data recently produced by Dexter Chua on $d_2$ differentials in
the Adams spectral sequence for the cofiber of $\eta$
shows that the classical element $g Q_3$ must be the target
of a hidden $\eta$ extension.  Therefore,
there is either a hidden $\eta$ extension from
$h_1^2 f_2$ to $\tau g Q_3$,
or from $\tau h_1 x_{85,6}$ to $\tau^2 g Q_3$.
\end{remark}
}

\begin{lemma}
\label{lem:eta-th1Q2}
\revdeg{58, 8, 30}
There is no hidden $\eta$ extension on $\tau h_1 Q_2$.
\end{lemma}

\begin{proof}
There cannot be a hidden $\eta$ extension from
$\tau h_1 Q_2$ to $\tau^2 \D h_1 d_0 g$ by comparison to $\tmf$.
It remains to show that there cannot be a hidden
$\eta$ extension from $\tau h_1 Q_2$ to $\tau M d_0$.

Note that $h_1 d_0 Q_2 = \tau^3 d_1 g^2$,
so $\kappa \{h_1 Q_2\}$ is detected by
$\tau^3 d_1 g^2$.  Therefore,
$\kappa \{h_1 Q_2\} + \tau \kappa_1 \kappabar^2$
is detected in higher filtration.
The only possibility is $\tau^3 e_0 g m$, but that cannot occur
by comparison to $\mmf$.
Therefore,
$\kappa \{h_1 Q_2\} + \tau \kappa_1 \kappabar^2$ is zero.

Now $\tau \eta \kappa_1 \kappabar^2$ is zero because
$\tau \eta \kappa_1 \kappabar$ cannot be detected by
$\D h_1 d_0^2$ by comparison to $\tmf$.
Therefore,
$\eta \kappa \{h_1 Q_2\}$ is zero,
so $\tau \eta \kappa \{h_1 Q_2\}$ is also zero.

On the other hand, $\tau \kappa \{M d_0\}$ is non-zero
because it is detected by $\tau M d_0^2$.
Therefore $\tau \eta \{h_1 Q_2\}$ cannot be detected
by $\tau M d_0$.
\end{proof}

\begin{lemma}
\label{lem:eta-th1^2h5^2}
\revdeg{64, 4, 33}
There is no hidden $\eta$ extension on $\tau h_1^2 h_5^2$.
\end{lemma}

\begin{proof}
Table \ref{tab:2-extn} shows that $\tau h_1^2 h_5^2$ is the
target of a hidden $2$ extension.
\end{proof}

\begin{lemma}
\label{lem:eta-t^2h1X2}
\revdeg{64, 8, 33}
There is a hidden $\eta$ extension from $\tau^2 h_1 X_2$
to $\tau^2 M h_0 g$.
\end{lemma}

\begin{proof}
Table \ref{tab:eta-extn} shows that there is a hidden
$\eta$ extension from $\tau h_1 X_2$ to $c_0 Q_2$.
Since $c_0 Q_2$ does not support a hidden $\tau$ extension,
there exists an element $\beta$ in $\pi_{65,35}$ that is 
detected by $c_0 Q_2$ such that $\tau \beta = 0$.

Projection from $C\tau$ to the top cell takes
$\ol{c_0 Q_2}$ and $P(A+A')$ to $c_0 Q_2$ and $\tau M h_0 h_2 g$
respectively.  Since $h_2 \cdot \ol{c_0 Q_2} = P(A+A')$
in the Adams spectral sequence for $C\tau$, it follows that
$\nu \beta$ is non-zero and detected by $\tau M h_0 h_2 g$.

Let $\alpha$ be an element of $\pi_{63,33}$ that is detected
by $\tau X_2 + \tau C'$, and
consider the sum $\eta^2 \alpha + \beta$.
Both terms are detected by $c_0 Q_2$, but the sum could be detected
in higher filtration.  In fact, the sum is non-zero because
$\nu (\eta^2 \alpha + \beta)$ is non-zero.

It follows that $\eta^2 \alpha + \beta$ is detected by $\tau M h_0 g$,
and that $\tau \eta^2 \alpha$ is detected by $\tau^2 M h_0 g$.
\end{proof}

\begin{lemma}
\label{lem:eta-th1^3h6}
\revdeg{66, 4, 34}
There is no hidden $\eta$ extension on $\tau h_1^3 h_6$.
\end{lemma}

\begin{proof}
The element $\tau \eta^2 \eta_6$ is detected by
$\tau h_1^3 h_6$.  Table \ref{tab:Toda} shows that
$\eta_6$ is contained in the Toda bracket
$\langle \eta, 2, \theta_5 \rangle$
Now shuffle to obtain
\[
\eta \cdot \tau \eta^2 \eta_6 = 4 \nu \eta_6 = 
4 \nu \langle \eta, 2, \theta_5 \rangle =
4 \langle \nu, \eta, 2 \rangle \theta_5,
\]
which equals zero because $2 \theta_5$ is zero.
\end{proof}

\begin{lemma}
\label{lem:eta-tD1h3^2}
\revdeg{66, 6, 35}
There is no hidden $\eta$ extension from
$\tau \D_1 h_3^2$ to $h_2^2 A'$.
\end{lemma}

\begin{proof}
Table \ref{tab:nu-extn} shows that $h_2^2 A'$
supports a hidden $\nu$ extension, so it cannot be the
target of a hidden $\eta$ extension.
\end{proof}

\begin{lemma}
\label{lem:eta-th1Q3}
\revdeg{68, 6, 36}
There is no hidden $\eta$ extension on $\tau h_1 Q_3$.
\end{lemma}

\begin{proof}
Table \ref{tab:nu-extn} shows that $\tau h_1 Q_3$
is the target of a hidden $\nu$ extension.
Therefore, it cannot be the source of a hidden $\eta$
extension.
\end{proof}

\begin{lemma}
\label{lem:eta-h3A'}
\revdeg{68, 7, 36}
There is a hidden $\eta$ extension from $h_3 A'$ to
$h_3 (\D e_1 + C_0)$.
\end{lemma}

\begin{proof}
Comparison to $C\tau$ shows that there is a hidden
$\eta$ extension from $h_3 A'$ to either
$\tau h_2^2 C' + h_3(\D e_1 + C_0)$ or
$h_3 (\D e_1 + C_0)$.
Table \ref{tab:nu-extn} shows that
$\tau h_2^2 C' + h_3(\D e_1 + C_0)$ supports a hidden
$\nu$ extension.  Therefore, it cannot be the target of a
hidden $\eta$ extension.
\end{proof}

\begin{lemma}
\label{lem:eta-h2^2h6}
\revdeg{69, 3, 36}
There is a hidden $\eta$ extension from $h_2^2 h_6$ to
$\tau h_0 h_2 Q_3$.
\end{lemma}

\begin{proof}
Table \ref{tab:Massey} gives the Massey product
$h_0 h_2 = \langle h_1, h_0, h_1 \rangle$.
Therefore,
\[
\langle \tau h_1 Q_3, h_0, h_1 \rangle = 
\{ \tau h_0 h_2 Q_3, \tau h_0 h_2 Q_3 + \tau h_1 h_3 H_1 \}.
\]
Table \ref{tab:nu-extn} shows that
there is a hidden $\nu$ extension from
$h_2 h_5^2$ to $\tau h_1 Q_3$, so
$\nu^2 \theta_5$ is detected by $\tau h_1 Q_3$.
Therefore,
the Toda bracket $\langle \nu^2 \theta_5, 2, \eta \rangle$
is detected by $\tau h_0 h_2 Q_3$ or by 
$\tau h_0 h_2 Q_3 + \tau h_1 h_3 H_1$.

Now $\langle \nu^2 \theta_5, 2, \eta \rangle$
contains $\nu^2 \langle \theta_5, 2, \eta \rangle$. This 
expression 
equals $\nu \theta_5 \langle 2, \eta, \nu \rangle$, which
equals zero because $\langle 2, \eta, \nu \rangle$ is contained
in $\pi_{5,3} = 0$.

We now know that $\langle \nu^2 \theta_5, 2, \eta \rangle$
equals its own determinacy, so
$\tau h_0 h_2 Q_3$ or $\tau h_0 h_2 Q_3 + \tau h_1 h_3 H_1$
detects a multiple of $\eta$.
The only possibility is that there is a hidden $\eta$
extension on $h_2^2 h_6$.

The target of this extension cannot be
$\tau h_0 h_2 Q_3 + \tau h_1 h_3 H_1$ by comparison 
to $C\tau$.
\end{proof}

\begin{lemma}
\label{lem:eta-h0^3h3h6}
\revdeg{70, 5, 36}
There is no hidden $\eta$ extension on $h_0^3 h_3 h_6$.
\end{lemma}

\begin{proof}
There are several possible targets for a hidden
$\eta$ extension on $h_0^3 h_3 h_6$.
The element $\tau \D^2 h_2 g$ is ruled out
because it supports an $h_2$ extension.
The element $\D^2 h_4 c_0$ is ruled out by comparison to $C\tau$.
The elements $\tau h_3^2 Q_2$ and $\tau d_0 Q_2$ are ruled out
because Table \ref{tab:eta-extn}
shows that they are targets of hidden $\eta$ extensions from
$\tau^2 h_1 h_3 H_1$ and $\tau^2 h_1 D'_3$ respectively.

The only remaining possibility is $\tau^2 l_1$.  This case is
more complicated.

Table \ref{tab:Toda} shows that $h_0^3 h_3 h_6$ detects the Toda bracket
$\langle 8\sigma, 2, \theta_5 \rangle$.
Now shuffle to obtain
\[
\eta \langle 8\sigma, 2, \theta_5 \rangle =
\langle \eta, 8 \sigma, 2 \rangle \theta_5.
\]
Table \ref{tab:Toda} shows that
$\langle \eta, 8 \sigma, 2 \rangle$ contains $\mu_9$
and has indeterminacy generated by
$\tau \eta^2 \sigma$ and $\tau \eta \epsilon$.
Thus the expression
$\langle \eta, 8 \sigma, 2 \rangle \theta_5$ contains
at most four elements.

The product $\mu_9 \theta_5$ is detected in filtration at least $8$,
so it is not detected by $\tau^2 l_1$.
The product $(\mu_9 + \tau \eta^2 \sigma) \theta_5$
is detected by $\tau h_1^2 p'$ because
Table \ref{tab:misc-extn}
shows that there is a hidden $\sigma$ extension
from $h_5^2$ to $p'$.
The product
$(\mu_9 + \tau \eta \epsilon) \theta_5$ is also 
detected by $\tau h_1^2 p' = \tau h_1 h_5^2 c_0$.
Finally, the product
$(\mu_9 + +\tau \eta^2 \sigma + \tau \eta \epsilon) \theta_5$
equals $(\mu_9 + \tau \nu^3) \theta_5$, which also
must be detected in filtration at least $8$. 
\end{proof}

\begin{lemma}
\label{lem:eta-h2Q3}
\revdeg{70, 6, 38}
There is no hidden $\eta$ extension on $h_2 Q_3$.
\end{lemma}

\begin{proof}
There cannot be a hidden $\eta$ extension on $\tau h_2 Q_3$
because it is a multiple of $h_2$.  Therefore, the possible
targets for an $\eta$ extension on $h_2 Q_3$ must be annihilated
by $\tau$.

The element $h_1^3 h_3 H_1$ cannot be the target because
Table \ref{tab:tau-extn} shows that it supports a hidden
$\tau$ extension.
The element $\tau M h_2^2 g$ cannot be the target because
Table \ref{tab:nu-extn} shows that it supports a hidden
$\nu$ extension to $M h_1 d_0^2$.
\end{proof}

\begin{lemma}
\label{lem:eta-th1h3H1}
\revdeg{70, 7, 37}
There is a hidden $\eta$ extension from $\tau h_1 h_3 H_1$ to
$h_3^2 Q_2$.
\end{lemma}

\begin{proof}
Table \ref{tab:eta-extn} shows that there is a hidden $\eta$ extension
from $\tau h_1 H_1$ to $h_3 Q_2$.  Now multiply by $h_3$.
\end{proof}

\begin{lemma}
\label{lem:eta-h1h3(De1+C0)}
\mbox{}
\begin{enumerate}
\item
\revdeg{70, 10, 38}
There is no hidden $\eta$ extension on $h_1 h_3 (\D e_1 + C_0)$.
\item
\revdeg{70, 10, 38}
There is no hidden $\eta$ extension on 
$\tau h_2 C'' + h_1 h_3 (\D e_1 + C_0)$.
\end{enumerate}
\end{lemma}

\begin{proof}
The element $\tau M h_2^2 g$ is the 
only possible target for such hidden $\eta$ extensions.
However, Table \ref{tab:nu-extn} shows that there is a hidden
$\nu$ extension from $\tau M h_2^2 g$ to $M h_1 d_0^2$.
\end{proof}

\begin{lemma}
\label{lem:eta-th1^2p'}
\revdeg{71, 6, 37}
There is no hidden $\eta$ extension on $\tau h_1^2 p'$.
\end{lemma}

\begin{proof}
The element $\tau h_1^2 p'$ detects 
$\tau \eta^2 \sigma \theta_5$ because Table \ref{tab:misc-extn}
shows that there is a hidden $\sigma$ extension from
$h_5^2$ to $p'$.
Then $\tau \eta^3 \sigma \theta_5$ is zero since
$\tau \eta^3 \sigma$ is zero.
\end{proof}

\begin{lemma}
\label{lem:eta-h1^2h3H1}
\revdeg{71, 8, 39}
There is no hidden $\eta$ extension on $h_2^3 H_1$.
\end{lemma}

\begin{proof}
Table \ref{tab:misc-extn} shows that $M d_0$ detects the product
$\kappa \theta_{4.5}$.  Then
Table \ref{tab:Toda} shows that $h_2^3 H_1$ detects the
Toda bracket $\langle \nu, \epsilon, \kappa \theta_{4.5} \rangle$.
Now shuffle to obtain
\[
\eta \langle \nu, \epsilon, \kappa \theta_{4.5} \rangle =
\langle \eta, \nu, \epsilon \rangle \kappa \theta_{4.5},
\]
which is zero because $\langle \eta, \nu, \epsilon \rangle$
is contained in $\pi_{13,8} = 0$.
\end{proof}

\begin{lemma}
\label{lem:eta-th1h6c0}
\revdeg{72, 5, 37}
There is a hidden $\eta$ extension from $\tau h_1 h_6 c_0$ 
to $\tau^2 h_2^2 Q_3$.
\end{lemma}

\begin{proof}
The hidden $\tau$ extension from $h_1^2 h_6 c_0$ to
$h_0 d_0 D_2$ implies that $\tau h_1 h_6 c_0$ must support a hidden
$\eta$ extension.  However, this hidden $\tau$ extension crosses
the $\tau$ extension from $\tau h_2^2 Q_3$ to $\tau^2 h_2^2 Q_3$.
Therefore, the target of the hidden $\eta$ extension is either
$\tau^2 h_2^2 Q_3$ or $h_0 d_0 D_2$.

The element $\tau h_1 h_6 c_0$ detects the product 
$\tau \eta_6 \epsilon$, so we want to compute 
$\tau \eta \eta_6 \epsilon$.
Table \ref{tab:Toda} shows that 
$\eta_6$ belongs to $\langle \theta_5, 2, \eta \rangle$.
Shuffle to obtain
\[
\tau \eta \eta_6 \epsilon = 
\langle \theta_5, 2, \eta \rangle \tau \eta \epsilon =
\theta_5 \langle 2, \eta, \tau \eta \epsilon \rangle.
\]
Table \ref{tab:Toda} shows that
$\langle 2, \eta, \tau \eta \epsilon \rangle$
contains $\zeta_{11}$.
Finally, $\theta_5 \zeta_{11}$ is detected by
$\tau^2 h_2^2 Q_3 = h_5^2 \cdot P h_2$.
\end{proof}

\begin{lemma}
\label{lem:eta-h1^3p'}
\revdeg{72, 7, 39}
There is no hidden $\eta$ extension on $h_1^3 p'$.
\end{lemma}

\begin{proof}
The element $h_1^3 p'$ does not support a hidden $\tau$ extension,
while Table \ref{tab:tau-extn} shows that
there is a hidden $\tau$ extension from $\tau h_2^2 C''$ to
$\D^2 h_1^2 h_4 c_0$.  Therefore, there cannot be a hidden
$\eta$ extension from $h_1^3 p'$ to $\tau h_2^2 C''$.
\end{proof}

\begin{lemma}
\label{lem:eta-h0d0D2}
\revdeg{72, 11, 38}
There is a hidden $\eta$ extension from $h_0 d_0 D_2$ to
$\tau M d_0^2$.
\end{lemma}

\begin{proof}
Table \ref{tab:Toda} shows that $h_0 d_0 D_2$ detects the
Toda bracket $\langle \tau \kappabar \theta_{4.5}, 2 \nu, \nu \rangle$.
Now shuffle to obtain
\[
\langle \tau \kappabar \theta_{4.5}, 2 \nu, \nu \rangle \eta =
\tau \kappabar \theta_{4.5} \langle 2 \nu, \nu, \eta \rangle.
\]
Table \ref{tab:Toda} shows that 
the Toda bracket $\langle 2 \nu, \nu, \eta \rangle$
contains $\epsilon$.
Finally, $\tau \kappabar \theta_{4.5} \epsilon$
is detected by $\tau M d_0^2$ because
Table \ref{tab:misc-extn} shows that there 
is a hidden $\epsilon$ extension from
$\tau M g$ to $M d_0^2$.
\end{proof}

\begin{lemma}
\label{lem:eta-h0h3d2}
\revdeg{75, 6, 40}
There is a hidden $\eta$ extension from $h_0 h_3 d_2$ 
to $\tau d_1 g_2$.
\end{lemma}

\begin{proof}
Table \ref{tab:Toda} shows that the Toda bracket
$\langle \eta, \sigma^2, \eta, \sigma^2 \rangle$
equals $\kappa_1$.
We would like to consider the shuffle
\[
\langle \eta, \sigma^2, \eta, \sigma^2 \rangle \tau \kappabar_2 =
\eta \langle \sigma^2, \eta, \sigma^2, \tau \kappabar_2 \rangle,
\]
but we must show that the Toda bracket
$\langle \eta, \sigma^2, \tau \kappabar_2 \rangle$ is well-defined and
contains zero.
It is well-defined because $\sigma^2 \kappabar_2$ is detected by
$h_3^2 g_2$ in $\pi_{58,32}$, and there are no $\tau$ extensions
on this group.
The bracket contains zero by comparison to $\tmf$, since
all non-zero elements of $\pi_{60,32}$ are detected
by $\tmf$.

We have now shown that $\tau \kappa_1 \kappabar_2$ is divisible
by $\eta$.  The only possibility is that 
there is a hidden $\eta$ extension from $h_0 h_3 d_2$ to
$\tau d_1 g_2$.
\end{proof}

\begin{lemma}
\label{lem:eta-tm1}
\mbox{}
\begin{enumerate}
\item
\revdeg{77, 3, 40}
There is no hidden $\eta$ extension on $h_3^2 h_6$.
\item
\revdeg{77, 7, 41}
There is no hidden $\eta$ extension on $\tau m_1$.
\end{enumerate}
\end{lemma}

\begin{proof}
Table \ref{tab:2-extn} shows that $e_0 A'$ and $\tau e_0 A'$ support hidden $2$ extensions, so they cannot be the targets of hidden
$\eta$ extensions.
\end{proof}

\begin{lemma}
\label{lem:eta-h0x77,7}
\revdeg{77, 8, 40}
There is no hidden $\eta$ extension on $h_0 x_{77,7}$.
\end{lemma}

\begin{proof}
Table \ref{tab:2-extn} shows that $h_0 x_{77,7}$ is the target
of a hidden $2$ extension.
\end{proof}

\begin{lemma}
\label{lem:eta-h1h6d0}
\revdeg{78, 6, 41}
There is no hidden $\eta$ extension on $h_1 h_6 d_0$.
\end{lemma}

\begin{proof}
Table \ref{tab:Toda} shows that $h_1 h_6$ detects the Toda bracket
$\langle \theta_5, 2, \eta \rangle$, so
$h_1 h_6 d_0$ detects
$\langle \theta_5, 2, \eta \rangle \kappa$.
Now shuffle to obtain
\[
\langle \theta_5, 2, \eta \rangle \eta \kappa = 
\theta_5 \langle 2, \eta, \eta \kappa \rangle.
\]
Table \ref{tab:Toda} shows that the Toda bracket
$\langle 2, \eta, \eta \kappa \rangle$ equals $\nu \kappa$.
Thus we need to compute the product
$\nu \kappa \theta_5$.
Lemma \ref{lem:sigma^2+kappa-h5^2} shows that this
product equals 
\rev{either $\nu \sigma^2 \theta_5$ or
$\nu (\sigma^2 \theta_5 + \tau^2 \kappa_1 \kappabar_2)$.
These expressions equal $0$ and $\tau^2 \nu \kappa_1 \kappabar_2$
respectively since $\nu \sigma = 0$.

It remains to show that $\tau^2 \nu \kappa_1 \kappabar_2$ is zero.
The proof of Lemma \ref{lem:eta-h0h3d2} shows that
$\tau \kappa_1 \kappabar_2$ is divisible by $\eta$.
Therefore, $\tau^2 \nu \kappa_1 \kappabar_2$ is zero
since $\eta \nu = 0$.
}
\end{proof}

\begin{lemma}
\label{lem:eta-th1^2x76,6}
\revdeg{78, 8, 41}
There is no hidden $\eta$ extension on $\tau h_1^2 x_{76,6}$.
\end{lemma}

\begin{proof}
Let $\alpha$ be an element of $\pi_{77,41}$ that is detected by
$h_1 x_{76,6}$.
Then $\tau h_1^2 x_{76,6}$ detects $\tau \eta \alpha$.
Now consider the shuffle
\[
\tau \eta^2 \alpha = \langle 2, \eta, 2 \rangle \alpha =
2 \langle \eta, 2, \alpha \rangle.
\]
Note that $2 \alpha$ is zero because there are no $2$ extensions
in $\pi_{77,41}$, so the second bracket is well-defined.

Finally, 
$2 \langle \eta, 2, \alpha \rangle$ must be zero because there are 
no $2$ extensions in $\pi_{79,42}$ in sufficiently high filtration.
\end{proof}

\rev{
\begin{lemma}
\label{lem:eta-h0^6h4h6}
\revdeg{78, 8, 40}
There is a hidden $\eta$ extension from
$h_0^6 h_4 h_6$ to $\tau \D B_6$.
\end{lemma}

\begin{proof}
In the homotopy of $C\tau$, there is a hidden $\eta$
extension from $h_0^6 h_4 h_6$ to $\D^2 n$.  Therefore,
$h_0^6 h_4 h_6$ must support a hidden $\eta$ extension
whose target lies in Adams filtration 13 or less.
However, $\D^2 n$ is not the target because it supports an
$h_2$ extension.  The only remaining possible
target is $\tau \D B_6$.
\end{proof}

}

\rev{
\begin{lemma}
\label{lem:eta-e0A'}
\revdeg{78, 10, 42}
There is a hidden $\eta$ extension from $e_0 A'$ to $\tau M e_0^2$.
\end{lemma}

\begin{proof}
The classical relation
$g C' = e_0 G_0$ implies the $\C$-motivic relation
$\tau g \cdot C' = e_0 \cdot \tau G_0$ modulo the possible
error term $\D j_1$.  The
error term does not appear because of $h_1^2$ extensions.

Table \ref{tab:Massey} shows that
$\tau G_0$ equals $\langle A', h_1, h_2 \rangle$.
Therefore, we have
\[
\tau g C' = e_0 \langle A', h_1, h_2 \rangle =
\langle e_0 A', h_1, h_2 \rangle.
\]
The second equality holds because there is no indeterminacy by 
inspection.

Let $\alpha$ be an element of $\pi_{78,44}$ that is detected by
$e_0 A'$.  If the product $\eta \alpha$ were zero, then the
Moss Convergence Theorem would imply that
$\tau g C'$ is a permanent cycle that detects the 
Toda bracket
$\langle \alpha, \eta, \nu \rangle$.
However, $\tau g C'$ supports a $d_4$ differential and does not survive.

We now know that $e_0 A'$ supports a hidden $\eta$ extension.
After ruling out $\tau^2 \D h_1 e_0^2 g$ by comparison to $\mmf$,
the only remaininng possible target is $\tau M e_0^2$.
\end{proof}
}

\begin{lemma}
\label{lem:eta-h2h4h6}
\revdeg{81, 3, 42}
If $h_2 h_4 h_6$ supports a hidden $\eta$ extension, then its
target is not $\tau h_2^2 x_{76,6}$.
\end{lemma}

\begin{proof}
Table \ref{tab:nu-extn} shows that
$\tau h_2^2 x_{76,6}$ supports a hidden $\nu$ extension, so
it cannot be the target of a hidden $\eta$ extension.
\end{proof}

\rev{
\begin{lemma}
\label{lem:eta-h1^3h4h6}
\revdeg{81, 5, 43}
There is no hidden $\eta$ extension on $h_1^3 h_4 h_6$.
\end{lemma}

\begin{proof}
The element $\tau h_1^3 h_4 h_6$ is a multiple of $h_0$, so
it cannot support a hidden $\eta$ extension.
This eliminates all possible targets except for $\tau (\D e_1 + C_0) g$.

However, 
$\tau (\D e_1 + C_0) g$ supports a hidden $\tau$ extension.
As in the previous paragraph, this eliminates
$\tau (\D e_1 + C_0) g$ as a possible target.
\end{proof}
}

\begin{lemma}
\label{lem:eta-h3^2n1}
\revdeg{81, 7, 44}
There is no hidden $\eta$ extension on $h_3^2 n_1$.
\end{lemma}

\begin{proof}
The element $\tau h_3^2 n_1 = h_3^2 (\tau Q_3 + \tau n_1)$
detects $\sigma^2 \{\tau Q_3 + \tau n_1\}$.  Then
$\eta \sigma^2 \{\tau Q_3 + \tau n_1\}$ is zero because
$\eta \sigma^2$ is zero.
\end{proof}

\begin{lemma}
\label{lem:eta-D^2p}
\revdeg{81, 12, 42}
There is no hidden $\eta$ extension on $\D^2 p$.
\end{lemma}

\begin{proof}
Table \ref{tab:nu-extn} shows that
$\D^2 p$ is the target of a hidden $\nu$ extension, so it cannot
be the source of an $\eta$ extension.
\end{proof}

\begin{lemma}
\label{lem:eta-h6c1}
\revdeg{82, 4, 43}
There is no hidden $\eta$ extension on $h_6 c_1$.
\end{lemma}

\begin{proof}
Table \ref{tab:Toda} shows that $h_6 c_1$ detects the Toda
bracket $\langle \sigmabar, 2, \theta_5 \rangle$.  By inspection,
all possible indeterminacy is in higher Adams filtration,
so $h_6 c_1$ detects every element of the Toda bracket.

Shuffle to obtain
\[
\eta \langle \sigmabar, 2, \theta_5 \rangle =
\langle \eta, \sigmabar, 2 \rangle \theta_5.
\]
The Toda bracket $\langle \eta, \sigmabar, 2 \rangle$
is detected in filtration at least $5$ since
the Massey product $\langle h_1, c_1, h_0 \rangle$
is zero.
Therefore, the Toda bracket equals $\{0, \eta \kappabar\}$.

We now know that $\eta \langle \sigmabar, 2, \theta_5 \rangle$
contains zero, and therefore $h_6 c_1$ does not support
a hidden $\eta$ extension.
\end{proof}

\begin{lemma}
\label{lem:eta-h0h6g}
\revdeg{83, 6, 44}
There is no hidden $\eta$ extension on $h_0 h_6 g$.
\end{lemma}

\begin{proof}
Table \ref{tab:Toda} shows that $h_0 h_6 g$ detects the Toda bracket
$\langle \nu, \eta, \eta_6 \kappa \rangle$.
Shuffle to obtain
\[
\eta \langle \nu, \eta, \eta_6 \kappa \rangle =
\langle \eta, \nu, \eta \rangle \eta_6 \kappa.
\]
Table \ref{tab:Toda} shows that $\langle \eta, \nu, \eta \rangle$
equals $\nu^2$.
Finally, 
\[
\nu^2 \eta_6 \kappa = \nu^2 \kappa \langle \eta, 2, \theta_5 \rangle =
\nu \theta_5 \kappa \langle \nu, \eta, 2 \rangle,
\]
which equals zero because $\langle \nu, \eta, 2 \rangle$ is contained
in $\pi_{5,3} = 0$.
\end{proof}

\begin{lemma}
\label{lem:eta-h2h4Q3}
\revdeg{85, 7, 46}
There is no hidden $\eta$ extension on $h_2 h_4 Q_3$.
\end{lemma}

\begin{proof}
We must eliminate $\tau h_2 g C'$ as a possible target.
One might hope to use the homotopy of $C\tau$ in order to do this,
but the homotopy of $C\tau$ has an $\eta$ extension in the
relevant degree that could possibly detect a hidden extension
from $h_2 h_4 Q_3$ to $\tau h_2 g C'$.


If there were a hidden $\eta$ extension from $h_2 h_4 Q_3$
to $\tau h_2 g C'$, then the hidden $\tau$ extension
from $\tau h_2 g C'$ to $\D^2 h_2^2 d_1$ would imply
that there is a hidden $\eta$ extension from
$\tau h_2 h_4 Q_3$ to $\D^2 h_2^2 d_1$.
However, $\tau h_2 h_4 Q_3$ detects the product
$\nu_4 \{\tau Q_3 + \tau n_1\}$, and $\eta \nu_4$ is zero.
Therefore,
$\tau h_2 h_4 Q_3$ cannot support a hidden $\eta$ extension.
\end{proof}

\begin{lemma}
\label{lem:eta-h1h6c0d0}
\mbox{}
\begin{enumerate}
\item
\revdeg{86, 9, 46}
There is no hidden $\eta$ extension on $h_1 h_6 c_0 d_0$.
\item
\revdeg{86, 10, 45}
There is no hidden $\eta$ extension on $P h_1 h_6 d_0$.
\end{enumerate}
\end{lemma}

\begin{proof}
Table \ref{tab:2-extn} shows that $h_1 h_6 c_0 d_0$ and
$P h_1 h_6 d_0$ are targets of hidden $2$ extensions, 
so they cannot be the sources of hidden $\eta$ extensions.
\end{proof}

\begin{lemma}
\label{lem:eta-h0^3h6i}
\revdeg{86, 11, 44}
There is a hidden $\eta$ extension from $h_0^3 h_6 i$ to
$\tau^2 \D^2 c_1 g$.
\end{lemma}

\begin{proof}
The Adams differential $d_2(\D^3 h_3^2) = \D^2 h_0^3 x$
implies that
$\tau^2 \D^2 c_1 g = h_1 \cdot \D^3 h_3^2$ detects the Toda bracket
$\langle \eta, 2, \{\D^2 h_0^2 x\} \rangle$.
However, the later Adams differential $d_5(h_0^2 h_6 i) = \D^2 h_0^2 x$
implies that $0$ belongs to $\{\D^2 h_0^2 x\}$.
Therefore,
$\tau^2 \D^2 c_1 g$ detects $\langle \eta, 2, 0 \rangle$,
so $\tau^2 \D^2 c_1 g$ detects a multiple of $\eta$.
The only possibility is that there is a hidden $\eta$ extension
from $h_0^3 h_6 i$ to $\tau^2 \D^2 c_1 g$.
\end{proof}

\begin{lemma}
\label{lem:eta-B6d1}
\revdeg{87, 11, 48}
There is no hidden $\eta$ extension on $B_6 d_1$.
\end{lemma}

\begin{proof}
Table \ref{tab:2-extn} shows that $B_6 d_1$ is the target of a 
hidden $2$ extension, so it cannot be the source of a hidden
$\eta$ extension.
\end{proof}

\begin{lemma}
\label{lem:eta-h1^2h4h6c0}
\revdeg{88, 7, 47}
There is no hidden $\eta$ extension on $h_1^2 h_4 h_6 c_0$.
\end{lemma}

\begin{proof}
Table \ref{tab:nu-extn} shows that $h_1^2 h_6 h_6 c_0$
is the target of a hidden $\nu$ extension, so it cannot support
a hidden $\eta$ extension.
\end{proof}

\begin{lemma}
\label{lem:eta-D^2h1f1}
\revdeg{89, 13, 47}
There is a hidden $\eta$ extension from $\D^2 h_1 f_1$ to
the element $\tau \D^2 h_2 c_1 g$.
\end{lemma}

\begin{proof}
The element $\tau \D^2 h_2 c_1 g$ detects the product
$\nu^2 \{\D^2 t\}$.
Table \ref{tab:Toda} shows that $\nu^2$ equals the Toda bracket
$\langle \eta, \nu, \eta \rangle$.
Shuffle to obtain
\[
\langle \eta, \nu, \eta \rangle \{\D^2 t\} =
\eta \langle \nu, \eta, \{\D^2 t\} \rangle.
\]
This shows that $\tau \D^2 h_2 c_1 g$ is the target of a hidden
$\eta$ extension.  The only possible source for this extension
is $\D^2 h_1 f_1$.
\end{proof}

\section{Hidden $\nu$ extensions}
\label{sctn:nu-extn}

\begin{thm}
\label{thm:nu-extn}
Tables \ref{tab:nu-extn} \rev{and \ref{tab:nu-extn-null}} list some hidden extensions by $\nu$.
\end{thm}

\begin{proof}
Many of the hidden extensions follow by comparison to 
$C\tau$.  For example, there is a hidden $\nu$ extension
from $ h_0^2 g$ to $h_1 c_0 d_0$ in the Adams spectral
sequence for $C\tau$.
Pulling back along inclusion of the bottom cell into $C\tau$,
there must also be a hidden $\nu$ extension from $h_0^2 g$
to $h_1 c_0 d_0$ in the Adams spectral sequence for the sphere.
This type of argument is indicated by the notation
$C\tau$ in the fourth column of Table \ref{tab:eta-extn}.

Next, Table \ref{tab:tau-extn} shows a hidden $\tau$
extension from $h_1 c_0 d_0$ to $P h_1 d_0$.
Therefore, there is also a hidden $\nu$ extension from
$\tau h_0^2 g$ to $P h_1 d_0$.
This type of argument is indicated by the notation $\tau$
in the fourth column of Table \ref{tab:eta-extn}.

Some extensions can be resolved by comparison to
$\tmf$ or to $\mmf$.
For example,
Table \ref{tab:unit-mmf} shows that 
the classical unit map $S \map \tmf$
takes $\{\D h_1 h_3\}$ in $\pi_{32}$ to a non-zero element
$\alpha$ of $\pi_{32} \tmf$ 
such that $\nu \alpha = \eta \kappa \kappabar$
in $\pi_{35} \tmf$.
Therefore, there must be a hidden $\nu$ extension
from $\D h_1 h_3$ to $\tau h_1 e_0^2$.

The proofs of several of the extensions in Table \ref{tab:nu-extn}
rely on analogous extensions in $\mmf$.  
Extensions in $\mmf$ have not been rigorously analyzed \cite{Isaksen18}.
However, the specific extensions from $\mmf$ that we need
are easily deduced from extensions in $\tmf$, together
with the multiplicative structure.
For example, there is a hidden $\nu$ extension in $\tmf$
from $\D h_1$ to $\tau d_0^2$.
Therefore, there is a hidden $\nu$ extension in $\mmf$
from $\D h_1 g$ to $\tau d_0^2 g$, and also a hidden
$\nu$ extension from $\tau \D h_1 g$ to $\tau^2 d_0 e_0^2$
in the homotopy groups of the sphere spectrum.
Note that $\mmf$ really is required here, since $d_0^2 g$ equals
zero in the homotopy of $\tmf$.

Many cases require more complicated arguments.  In stems up to
approximately dimension 62, see 
\cite{Isaksen14c}*{Section 4.2.4 and Tables 31--32} and
\cite{WangXu18}.
The higher-dimensional cases are handled in the following lemmas.
\end{proof}

\begin{remark}
\label{rem:nu-extn-notes}
The last column of Table \ref{tab:nu-extn} indicates which
$\nu$ extensions are crossing, as well as which extensions
have indeterminacy in the sense of Section \ref{subsctn:hid-indet}.
\end{remark}

\begin{remark}
\label{rem:nu-h2h5d0}
The hidden $\nu$ extension from $h_2 h_5 d_0$ to $\tau g n$
is proved in \cite{WangXu18}, which relies on the
``$\R P^\infty$-method" to establish a hidden
$\sigma$ extension from $\tau h_3 d_1$ to $\D h_2 c_1$
and a hidden $\eta$ extension from $\tau h_1 g_2$ to $\D h_2 c_1$.
We now have easier proofs for these $\eta$ and $\sigma$ extensions,
using the hidden $\tau$ extension from $h_1^2 g_2$ to $\D h_2 c_1$
given in Table \ref{tab:tau-extn}, as well as the relation
$h_3^2 d_1 = h_1^2 g_2$.
\end{remark}

\begin{remark}
\label{rem:nu-t(De1+C0)g}
If $M \D h_1^2 d_0$ is not hit by a differential, then there is a 
hidden $\tau$ extension from $\tau M h_0 g^2$ from $M \D h_1^2 d_0$.
This implies that there must be a hidden
$\nu$ extension from $\tau (\D e_1 + C_0) g$ to $M \D h_1^2 d_0$.
\end{remark}

\begin{thm}
\label{thm:nu-extn-possible}
Table \ref{tab:nu-extn-possible}
lists all unknown hidden $\nu$ extensions, through
the $90$-stem. 
\end{thm}

\begin{proof}
Many possible extensions can be eliminated by
comparison to $C\tau$, to $\tmf$, or to $\mmf$.
For example,
there cannot be a hidden $\nu$ extension from $h_0 h_2 h_4$
to $\tau h_1 g$ by comparison to $C\tau$.

Other possibilities are eliminated by consideration of other
parts of the multiplicative structure.  For example,
there cannot be a hidden $\nu$ extension whose target supports
a multiplication by $\eta$, since $\eta \nu$ equals zero.

Many cases are eliminated by more complicated arguments.
These are handled in the following lemmas.
\end{proof}

\begin{remark}
\label{rem:h2-th1p1}
Comparison to synthetic homotopy eliminates several
possible hidden $\nu$ extensions, including:
\begin{enumerate}
\item
from $\tau h_1 p_1$ to $\tau x_{74,8}$.
\item
from $\D^2 p$ to $\tau M \D h_1 d_0$.
\end{enumerate}
See \cite{BIX} for more details.
\end{remark}

\begin{remark}
\label{rem:h0h2h4h6}
If $M \D h_1^2 d_0$ is not hit by a differential, then
$M \D h_1 d_0$ supports an $h_1$ extension, and
there cannot be a hidden $\nu$ extension from
$h_0 h_2 h_4 h_6$ to $M \D h_1 d_0$.
\end{remark}

\begin{lemma}
\label{lem:nu-De1+C0}
\revdeg{62, 8, 33}
There is a hidden $\nu$ extension from $\D e_1 + C_0$ to
$\tau M h_0 g$.
\end{lemma}

\begin{proof}
Table \ref{tab:Toda} shows that
$2 \kappabar$ is contained in 
$\tau \langle \nu, \eta, \eta \kappa \rangle$.
Shuffle to obtain that
\[
\nu \langle \eta, \eta \kappa, \tau \theta_{4.5} \rangle =
\langle \nu, \eta, \eta \kappa \rangle \tau \theta_{4.5},
\]
so $2 \kappabar \theta_{4.5}$ is divisible by $\nu$.

Table \ref{tab:misc-extn} shows that 
$\tau M g$ detects $\kappabar \theta_{4.5}$, so
$\tau M h_0 g$ detects 
$2 \kappabar \theta_{4.5}$.
Now we know that there is a hidden $\nu$ extension whose target is
$\tau M h_0 g$, and the only possible source is $\D e_1 + C_0$.
\end{proof}

\begin{remark}
\label{rem:eta,etakappa,ttheta4.5}
One consequence of the proof of Lemma \ref{lem:nu-De1+C0}
is that $\D e_1 + C_0$ detects the Toda bracket
$\langle \eta, \eta \kappa, \tau \theta_{4.5} \rangle$.
\end{remark}

\begin{lemma}
\label{lem:nu-th1H1}
\revdeg{63, 6, 33}
There is a hidden $\nu$ extension from $\tau h_1 H_1$ to
$\tau^2 M h_1 g$.
\end{lemma}

\begin{proof}
Lemma \ref{lem:kappa,2,eta} shows that 
the bracket $\langle \kappa, 2, \eta \rangle$
contains zero with indeterminacy generated by $\eta \rho_{15}$.
The bracket $\langle \tau \eta \theta_{4.5}, \kappa, 2 \rangle$
equals zero since $\pi_{61,32}$ is zero.
Therefore, the Toda bracket 
$\langle \tau \eta \theta_{4.5}, \kappa, 2, \eta \rangle$ is well-defined.

Table \ref{tab:Toda} shows that $\tau g $ detects 
$\langle \kappa, 2, \eta, \nu \rangle$.
Therefore, $\tau^2 M h_1 g$ detects
\[
\tau \eta \theta_{4.5} \langle \kappa, 2, \eta, \nu \rangle =
\langle \tau \eta \theta_{4.5}, \kappa, 2, \eta \rangle \nu.
\]
This shows that
$\tau^2 M h_1 g$ is the target of a $\nu$ extension,
and the only possible source is $\tau h_1 H_1$.
\end{proof}

\begin{remark}
\label{rem:tetatheta4.5,kappa,2,eta}
The proof of Lemma \ref{lem:nu-th1H1} shows that
$\tau h_1 H_1$ detects 
the Toda bracket
$\langle \tau \eta \theta_{4.5}, \kappa, 2, \eta \rangle$.
\end{remark}

\begin{lemma}
\label{lem:nu-h1h6}
\revdeg{64, 2, 33}
There is no hidden $\nu$ extension on $h_1 h_6$.
\end{lemma}

\begin{proof}
Table \ref{tab:Toda} shows that $h_1 h_6$ detects the Toda bracket
$\langle \eta, 2, \theta_5 \rangle$.  Shuffle to obtain
\[
\nu \langle \eta, 2, \theta_5 \rangle =
\langle \nu, \eta, 2 \rangle \theta_5 = 0,
\]
since $\langle \nu, \eta, 2 \rangle$ is contained in $\pi_{5,3} = 0$.
\end{proof}

\begin{lemma}
\label{lem:nu-h3Q2}
\revdeg{64, 8, 34}
There is no hidden $\nu$ extension on $h_3 Q_2$.
\end{lemma}

\begin{proof}
Table \ref{tab:eta-extn} shows that 
$\tau^2 \D h_2^2 e_0 g$ supports a hidden $\eta$ extension.
Therefore, it cannot be the target of a $\nu$ extension.
\end{proof}

\begin{lemma}
\label{lem:nu-h2^2A'}
\revdeg{67, 8, 36}
There is a hidden $\nu$ extension from
the element $h_2^2 A'$ to $h_1 h_3 (\D e_1 + C_0)$.
\end{lemma}

\begin{proof}
By comparison to $C\tau$, 
There cannot be a hidden $\nu$ extension from
$h_2^2 A'$ to $\tau h_2 C'' + h_1 h_3 (\D e_1 + C_0)$


Table \ref{tab:Toda} shows that
$\D e_1 + C_0$ detects the Toda bracket
$\langle \eta, \eta \kappa, \tau \theta_{4.5} \rangle$,
and $h_2 A'$ detects the Toda bracket
$\langle \nu, \eta, \tau \kappa \theta_{4.5} \rangle$.
Note that $h_2 A'$ also detects
$\langle \nu, \eta \kappa, \tau \theta_{4.5} \rangle$.

Now shuffle to obtain
\[
(\eta \sigma + \epsilon) \langle \eta, \eta \kappa, 
\tau \theta_{4.5} \rangle +
\nu^2 \langle \nu, \eta \kappa, \tau \theta_{4.5} \rangle =
\left\langle
\left[
\begin{array}{cc}
\eta \sigma + \epsilon & \nu^2
\end{array}
\right],
\left[
\begin{array}{c}
\eta \\
\nu 
\end{array}
\right],
\eta \kappa \right\rangle \tau \theta_{4.5}.
\]
The matric Toda bracket
$\displaystyle
\left\langle
\left[
\begin{array}{cc}
\eta \sigma + \epsilon & \nu^2
\end{array}
\right],
\left[
\begin{array}{c}
\eta \\
\nu 
\end{array}
\right],
\eta \kappa \right\rangle$
must equal $\{ 0, \nu^2 \sigmabar \}$,
since $\nu^2 \sigmabar = \{ h_1^2 h_4 c_0 \}$
is the only non-zero element of $\pi_{25,15}$,
and that element belongs to the indeterminacy because it is a
multiple of $\nu^2$.

Next observe that $\tau \nu^2 \sigmabar \theta_{4.5}$ is zero
because all possible values of
$\sigmabar \theta_{4.5}$ are multiples of $\eta$.
This shows that
\[
(\eta \sigma + \epsilon) \alpha + \nu^2 \beta = 0,
\]
for some $\alpha$ and $\beta$ detected by $\D e_1 + C_0$
and $h_2 A'$ respectively.
The product
$(\eta \sigma + \epsilon) \alpha$ is detected by
$h_1 h_3 (\D e_1 + C_0)$, so there must be a hidden $\nu$
extension from $h_2^2 A'$ to $h_1 h_3 (\D e_1 + C_0)$.
\end{proof}

\begin{lemma}
\label{lem:nu-h3A'}
\revdeg{68, 7, 36}
There is no hidden $\nu$ extension on $h_3 A'$.
\end{lemma}

\begin{proof}
Table \ref{tab:Toda} shows that $h_3 A'$ detects the Toda bracket
$\langle \sigma, \kappa, \tau \eta \theta_{4.5} \rangle$.
Now shuffle to obtain
\[
\nu \langle \sigma, \kappa, \tau \eta \theta_{4.5} \rangle = 
\langle \nu, \sigma, \kappa \rangle \tau \eta \theta_{4.5} =
\langle \eta, \nu, \sigma \rangle \tau \kappa \theta_{4.5}.
\]
The Toda bracket
$\langle \eta, \nu, \sigma \rangle$ is zero because
it is contained in $\pi_{12,7} = 0$.
\end{proof}

\begin{lemma}
\label{lem:nu-p'}
\revdeg{69, 4, 36}
There is no hidden $\nu$ extension on $p'$.
\end{lemma}

\begin{proof}
Table \ref{tab:misc-extn} shows that 
$p'$ detects the product $\sigma \theta_5$.
Therefore, it cannot support a hidden $\nu$ extension.
\end{proof}

\begin{lemma}
\label{lem:nu-h2^2C'}
\revdeg{69, 9, 38}
There is a hidden $\nu$ extension from $h_2^2 C'$ to $\tau^2 d_1 g^2$.
\end{lemma}

\begin{proof}
Let $\alpha$ be an element of $\pi_{63,33}$ that is detected by
$\tau X_2 + \tau C'$.
Table \ref{tab:misc-extn} shows that 
$\epsilon \alpha$ is detected by $d_0 Q_2$, so
$\eta \epsilon \alpha$ is detected by $\tau^3 d_1 g^2$.
On the other hand,
$\eta \sigma \alpha$ is zero by comparison to $C\tau$.

Now consider the relation $\eta^2 \sigma + \nu^3 = \eta \epsilon$.
This shows that $\nu^3 \alpha$ is detected by $\tau^3 d_1 g^2$.
Since $\nu^2 \alpha$ is detected by $\tau h_2^2 C'$,
there must be a hidden $\nu$ extension from
$h_2^2 C'$ to $\tau^2 d_1 g^2$.
\end{proof}

\begin{lemma}
\label{lem:nu-th1D3'}
\revdeg{70, 9, 37}
There is a hidden $\nu$ extension from $\tau h_1 D'_3$ to
$\tau M d_0^2$.
\end{lemma}

\begin{proof}
Table \ref{tab:Toda} shows that $\tau h_1 D'_3$ detects the 
Toda bracket $\langle \eta, \nu, \tau \kappabar \theta_{4.5} \rangle$.
Now shuffle to obtain
\[
\nu \langle \eta, \nu, \tau \kappabar \theta_{4.5} \rangle = 
\langle \nu, \eta, \nu \rangle \tau \kappabar \theta_{4.5}.
\]
The bracket $\langle \nu, \eta, \nu \rangle$
equals $\eta \sigma + \epsilon$ \cite{Toda62}.

Now we must compute 
$(\eta \sigma + \epsilon) \tau \kappabar \theta_{4.5}$.
The product $\sigma \kappabar$ is zero, and
Table \ref{tab:misc-extn} shows that
$\epsilon \kappabar \theta_{4.5}$ is detected by $M d_0^2$.
These two observations imply that
$(\eta \sigma + \epsilon) \tau \kappabar \theta_{4.5}$
is detected by $\tau M d_0^2$.
\end{proof}

\begin{lemma}
\label{lem:nu-h6c0}
\revdeg{71, 4, 37}
There is no hidden $\nu$ extension on $h_6 c_0$.
\end{lemma}

\begin{proof}
Table \ref{tab:Toda} shows that $h_6 c_0$ detects the Toda bracket
$\langle \epsilon, 2, \theta_5 \rangle$.  Now shuffle to obtain
\[
\nu \langle \epsilon, 2, \theta_5 \rangle = 
\langle \nu, \epsilon, 2 \rangle \theta_5.
\]
Finally, the Toda bracket $\langle \nu, \epsilon, 2 \rangle$ is 
zero because it is contained in $\pi_{12,7} = 0$.
\end{proof}

\begin{lemma}
\label{lem:nu-h2^2C''}
\revdeg{73, 11, 41} 
There is a hidden $\nu$ extension from $h_2^2 C''$ to $\tau g^2 t$.
\end{lemma}

\begin{proof}
Let $\alpha$ be an element of $\pi_{53,30}$ that is detected by $i_1$.
Table \ref{tab:nu-extn} shows $g t$ detects $\nu \alpha$.
Therefore $\tau g^2 t$ detects $\nu \kappa \alpha$,
so $\tau g^2 t$ must be the target of a hidden $\nu$ extension.
The element $h_2^2 C''$ is the only
possible source for this extension.
\end{proof}

\begin{lemma}
\label{lem:nu-Mh1h3g}
\revdeg{73, 12, 41}
There is no hidden $\nu$ extension on $M h_1 h_3 g$.
\end{lemma}

\begin{proof}
If there were a hidden $\nu$ extension from $M h_1 h_3 g$ to
$\tau g^2 t$, then there would also be a hidden $\nu$ extension
with target $\tau^2 g^2 t$.  But there is no possible source
for such an extension.
\end{proof}

\begin{lemma}
\label{lem:nu-h0h3d2}
\revdeg{75, 6, 40}
If there is a hidden $\nu$ extension on $h_0 h_3 d_2$, then
its target is $M \D h_1^2 h_3$.
\end{lemma}

\begin{proof}
The only other possible target is $e_0 A'$.
However, Table \ref{tab:2-extn} shows that $e_0 A'$ supports
a hidden $2$ extension, while $h_0 h_3 d_2$ does not.
\end{proof}

\begin{lemma}
\label{lem:nu-td1g2}
\revdeg{76, 8, 41}
There is no hidden $\nu$ extension on $\tau d_1 g_2$.
\end{lemma}

\begin{proof}
Table \ref{tab:eta-extn} shows that $\tau d_1 g_2$ is the target
of a hidden $\eta$ extension.  Therefore, it cannot be the source
of a hidden $\nu$ extension.
\end{proof}

\begin{lemma}
\label{lem:nu-h0h4A}
\revdeg{76, 8, 40}
There is no hidden $\nu$ extension on $h_0 h_4 A$.
\end{lemma}

\begin{proof}
\rev{
Table \ref{tab:misc-extn} shows that
$h_0 h_4 A$ detects either $\sigma^2 \theta_5$
or $\sigma^2 \theta_5 + \tau^2 \kappa_1 \kappabar_2$.
As in the proof of Lemma \ref{lem:eta-h1h6d0}, both possibilities
are annihilated by $\nu$.
}
\end{proof}

\begin{lemma}
\label{lem:nu-h3^2h6}
\revdeg{77, 3, 40}
There is a hidden $\nu$ extension from $h_3^2 h_6$ to $\tau h_1 x_1$.
\end{lemma}

\begin{proof}
Table \ref{tab:Toda} shows that
$h_3^2 h_6$ detects
$\langle \theta_5, 2, \sigma^2 \rangle$.
Let $\alpha$ be an element of $\pi_{77,40}$ that is contained
in this Toda bracket.
Then $\nu \alpha$ is an element of
\[
\langle \theta_5, 2, \sigma^2 \rangle \nu =
\langle \theta_5, 2 \sigma, \sigma \rangle \nu =
\theta_5 \langle 2 \sigma, \sigma, \nu \rangle \subseteq
\langle 2 \sigma, \sigma \theta_5, \nu \rangle.
\]

Table \ref{tab:misc-extn} shows that $p'$ detects 
$\sigma \theta_5$.
Therefore, the Toda bracket
$\langle 2 \sigma, \sigma \theta_5, \nu \rangle$
is detected by an element of the Massey product
$\langle h_0 h_3, p', h_2 \rangle$.
Table \ref{tab:Massey} shows that $h_0 e_2$ equals the
Massey product $\langle h_3, p', h_2\rangle$. 
By inspection of indeterminacy,
the Massey product $\langle h_0 h_3, p', h_2 \rangle$
contains $h_0^2 e_2 = \tau h_1 x_1$ with indeterminacy generated by
$h_0 h_6 e_0$.

We have now shown that 
$\nu \alpha$ is detected by either
$\tau h_1 x_1$ or $\tau h_1 x_1 + h_0 h_6 e_0$.
But $h_0 h_6 e_0 = h_2 h_6 d_0$ is a multiple of $h_2$,
so we may add an element in higher Adams filtration to $\alpha$,
if necessary, to conclude that $\nu \alpha$ is detected by
$\tau h_1 x_1$. 
\end{proof}

\begin{lemma}
\label{lem:nu-h0^7h4h6}
\revdeg{78, 9, 40}
There is a hidden $\nu$ extension from the element $h_0^7 h_4 h_6$ to
$\tau \D^2 h_1 d_1$.
\end{lemma}

\begin{proof}
Table \ref{tab:nu-extn} shows that there is a hidden
$\nu$ extension from $h_0^6 h_4 h_6$ to $\D^2 p$.
Therefore, there is also a hidden $\nu$ extension from
$h_0^7 h_4 h_6$ to $h_0 \cdot \D^2 p = \tau \D^2 h_1 d_1$.
\end{proof}

\begin{lemma}
\label{lem:nu-e0A'}
\revdeg{78, 10, 42}
There is no hidden $\nu$ extension on $e_0 A'$.
\end{lemma}

\begin{proof}
A possible hidden $\nu$ extension from $e_0 A'$ to $\D h_1^2 B_6$
would be detected by $C\tau$, but we have to be careful with the
analysis of the homotopy of $C\tau$ because of the $h_2$ extension
from $\ol{\D h_1 d_1 g}$ to $\D h_1^2 B_6$ in the Adams
$E_\infty$-page for $C\tau$.

Let $\alpha$ be an element of $\pi_{75,40} C\tau$ that is detected
by $\ol{h_3 C'}$.  Then $\nu \alpha$ is detected by
$e_0 A'$, and $\nu \alpha$ maps to zero under projection to the 
top cell because $h_3 C'$ does not support a $\nu$ extension
in the homotopy of the sphere.

Therefore, $\nu \alpha$ lies in the image of 
$e_0 A'$ under inclusion of the bottom
cell.  Since $\nu^2 \alpha$ is zero, $e_0 A'$ cannot support
a hidden $\nu$ extension to $\D h_1^2 B_6$.
\end{proof}

\begin{lemma}
\label{lem:nu-h1h4h6}
\revdeg{79, 3, 41}
There is no hidden $\nu$ extension on $h_1 h_4 h_6$.
\end{lemma}

\begin{proof}
Table \ref{tab:Toda} shows that $h_1 h_4 h_6$
detects the Toda bracket $\langle \eta_4, 2, \theta_5 \rangle$.
Shuffle to obtain
\[
\nu \langle \eta_4, 2, \theta_5 \rangle =
\langle \nu, \eta_4, 2 \rangle \theta_5.
\]
Finally, $\langle \nu, \eta_4, 2 \rangle$ must contain zero
in $\pi_{20,11}$ because $\tmf$ detects every element of
$\pi_{20,11}$.
\end{proof}

\begin{lemma}
\label{lem:nu-h3^2n1}
\revdeg{81, 7, 44}
There is no hidden $\nu$ extension on $h_3^2 n_1$.
\end{lemma}

\begin{proof}
The element $h_2 g D_3$ cannot be the target of a hidden
$\nu$ extension by comparison to $C\tau$.

The element $\tau h_3^2 n_1 = h_3 \cdot h_3 (\tau Q_3 + \tau n_1)$
detects a multiple of $\sigma$, so it cannot support a hidden
$\nu$ extension.  This rules out $h_2 g A'$ as a possible target.
\end{proof}

\begin{lemma}
\label{lem:nu-te1g2}
\revdeg{82, 8, 44}
There is no hidden $\nu$ extension on $\tau e_1 g_2$.
\end{lemma}

\begin{proof}
After eliminating other possibilities by comparison to
$\tmf$, comparison to $\mmf$, and by inspection of $h_1$
multiplications, the only possible target for a hidden
$\nu$ extension is $P h_1 x_{76,6}$.  

Let $\alpha$ be an element of $\pi_{82,45}$ that is detected
by $e_1 g_2$.  Then $\nu \alpha$ is detected by
$h_2 e_1 g_2 = h_1^3 h_4 Q_3$.
Choose an element $\beta$ of $\pi_{83,45}$ that is detected
by $h_1 h_4 Q_3$ such that $\tau \beta$ is zero.
Then $\eta^2 \beta$ is also detected by $h_1^3 h_4 Q_3$.
However, $\nu \alpha + \eta^2 \beta$ is not necessarily zero; 
it could be detected in Adams filtration at least $13$.
In any case,
$\tau \nu \alpha$ equals $\tau \eta^2 \beta = 0$ modulo
filtration $13$.
In particular, $\tau \nu \alpha$ cannot be detected by
$P h_1 x_{76,6}$ in filtration $11$.
\end{proof}

\begin{lemma}
\label{lem:nu-Ph0h2h6}
\revdeg{82, 11, 42}
There is no hidden $\nu$ extension on $P^2 h_0 h_2 h_6$.
\end{lemma}

\begin{proof}
Table \ref{tab:nu-extn} shows that there is a hidden
$\nu$ extension from $P^2 h_2 h_6$ to $\D^2 h_0 x$.
The target of a hidden $\nu$ extension
on $P^2 h_0 h_2 h_6$ must have Adams filtration greater than
the filtration of $\D^2 h_0 x$.
The only possibilities are ruled out by comparison to $\tmf$.
\end{proof}

\begin{lemma}
\label{lem:nu-(De1+C0)g}
\revdeg{82, 12, 45}
There is a hidden $\nu$ extension
from $(\D e_1 + C_0) g$ to $\tau M h_0 g^2$.
\end{lemma}

\begin{proof}
Let $\alpha$ be an element of $\pi_{62,33}$ that is detected
by $\D e_1 + C_0$.
Table \ref{tab:Toda} shows that
$(\D e_1 + C_0) g$ detects $\langle \alpha, \eta^3, \eta_4 \rangle$.
Then
\[
\nu \langle \alpha, \eta^3, \eta_4 \rangle =
\langle \nu \alpha, \eta^3, \eta_4 \rangle
\]
by inspection of indeterminacies.
Table \ref{tab:nu-extn} shows that $\tau M h_0 g$ detects $\nu \alpha$.
The Toda bracket
$\langle \nu \alpha, \eta^3, \eta_4 \rangle$ is detected by
the Massey product
\[
\langle \tau M h_0 g, h_1^3, h_1 h_4 \rangle =
\langle \tau M h_0 g, h_1^4, h_4 \rangle = 
M h_0 g \langle \tau, h_1^4, h_4 \rangle =
\tau M h_0 g^2.
\]
\end{proof}

\begin{lemma}
\label{lem:nu-h2c1A'}
\revdeg{83, 10, 45}
There is no hidden $\nu$ extension on $h_2 c_1 A'$.
\end{lemma}

\begin{proof}
Table \ref{tab:Toda} shows that 
$\tau h_2 c_1 A'$ detects 
$\langle \tau \theta_{4.5} \kappa, \eta, \nu \rangle \tau \sigmabar$.
Shuffle to obtain
\[
\langle \tau \theta_{4.5} \kappa, \eta, \nu \rangle \tau \sigmabar \nu =
\tau \theta_{4.5} \kappa \langle \eta, \nu, \tau \nu \sigmabar \rangle.
\]
The Toda bracket $\langle \eta, \nu, \tau \nu \sigmabar \rangle$
is zero because $\pi_{27, 15}$ contains only a $v_1$-periodic element
detected by $P^3 h_1^3$.

We now know that $\tau h_2 c_1 A'$ does not support a hidden
$\nu$ extension.  In particular, there cannot be a hidden
$\nu$ extension from $\tau h_2 c_1 A'$ to $M \D h_0^2 e_0$.
The hidden $\tau$ extension from $\tau^2 M h_1 g^2$ to 
$M \D h_0^2 e_0$ implies that there cannot be a hidden
$\nu$ extension from $h_2 c_1 A'$ to $\tau^2 M h_1 g^2$.

Additional cases are ruled out by comparison to $C\tau$ and
to $\mmf$.
\end{proof}

\begin{lemma}
\label{lem:nu-Dj1+tgC'}
\mbox{}
\begin{enumerate}
\item
\revdeg{83, 11, 45}
There is a hidden $\nu$ extension from $\D j_1 + \tau g C'$ to
$\tau^2 M h_1 g^2$.
\item
\revdeg{83, 11, 44}
There is a hidden $\nu$ extension from $\tau^2 g C'$ to $M \D h_0^2 e_0$.
\end{enumerate}
\end{lemma}

\begin{proof}
Table \ref{tab:nu-extn} shows that there exists an element
$\alpha$ in $\pi_{63,33}$ detected by $\tau h_1 H_1$ such that
$\nu$ is detected by $\tau^2 M h_1 g$.
(Beware that there is a crossing extension here, so not every
element detected by $\tau h_1 H_1$ has the desired property.)
Table \ref{tab:misc-extn} shows that 
$\tau^2 M h_1 g$ also detects $\tau \theta_{4.5} \eta \kappabar$.
However, $\nu \alpha$ does not necessarily equal
$\tau \theta_{4.5} \eta \kappabar$ because the difference could
be detected in higher filtration by $\D^2 h_1^3 h_4$.
In any case,
$\nu \kappabar \alpha$ equals $\tau \theta_{4.5} \eta \kappabar^2$.

The product $\theta_{4.5} \eta \kappabar^2$
is detected by $\tau^2 M h_1 g^2$.
The hidden $\tau$ extension from $\tau^2 M h_1 g^2$ to
$M \D h_0^2 e_0$ then implies that
$\nu \kappabar \alpha = \tau \theta_{4.5} \eta \kappabar^2$
is detected by $M \D h_0^2 e_0$.

We now know that $M \D h_0^2 e_0$ is the target of a hidden
$\nu$ extension.
The only possible source is $\tau^2 g C'$. (Lemma \ref{lem:nu-h2c1A'}
eliminates another possible source.)
This establishes the second extension.
The first extension follows from onsideration of $\tau$ extensions.
\end{proof}

\begin{remark}
\label{rem:nu-Dj1+tgC'}
The proof of Lemma \ref{lem:nu-Dj1+tgC'} shows that
$\nu \kappabar \alpha$ is detected by $M \D h_0^2 e_0$,
where $\alpha$ is detected by $\tau h_1 H_1$.
Note that $\kappabar \alpha$ is detected by $\tau^2 h_1 H_1 g =
\tau h_2 c_1 A'$.  But this does not show that 
$\tau h_2 c_1 A'$ supports a hidden $\nu$ extension.  Rather,
it shows that the source of the hidden $\nu$ extension is
either $\tau h_2 c_1 A'$, or a non-zero element in higher filtration.
\end{remark}

\begin{lemma}
\label{lem:nu-h2^2h4h6}
\revdeg{84, 4, 44}
There is no hidden $\nu$ extension on $h_2^2 h_4 h_6$.
\end{lemma}

\begin{proof}
Table \ref{tab:Toda} shows that $h_2^2 h_4 h_6$ detects the 
Toda bracket $\langle \nu \nu_4, 2, \theta_5 \rangle$.
Shuffle to obtain
\[
\nu \langle \nu \nu_4, 2, \theta_5 \rangle = 
\langle \nu, \nu \nu_4, 2 \rangle \theta_5.
\]
The Toda bracket $\langle \nu, \nu \nu_4, 2 \rangle$
is zero because $\pi_{25,14}$ consists only of a $v_1$-periodic
element detected by $P^2 h_1 c_0$.
\end{proof}

\rev{
\begin{lemma}
\label{lem:nu-tx85,6+h0^3c3}
\revdeg{85, 6, 44}
If $\tau x_{85,6} + h_0^3 c_3$ survives, then
it supports a hidden $\nu$ extension to
$h_1 x_{87,7} + \tau^2 g_2^2$.
\end{lemma}

\begin{proof}
By comparison to $C\tau$, there must be a hidden $\nu$ extension
whose target is either $h_1 x_{87,7}$ or $h_1 x_{87,7} + \tau^2 g_2^2$.

Table \ref{tab:nu-extn} shows that there is a hidden
$\nu$ extension from $\tau^2 h_2 h_4 Q_3$ to $\tau^2 h_0 g_2^2$.
This implies that the target of the $\nu$ extension on
$\tau x_{85,6} + h_0^3 c_3$ must be $h_1 x_{87,7} + \tau^2 g_2^2$.
\end{proof}
}

\begin{lemma}
\label{lem:nu-P^2h6c0}
\revdeg{87, 12, 45}
There is no hidden $\nu$ extension on $P^2 h_6 c_0$.
\end{lemma}

\begin{proof}
Table \ref{tab:misc-extn} shows that $P^2 h_6 c_0$ 
detects the product $\rho_{23} \eta_6$, and
$\nu \rho_{23} \eta_6$ is zero.
\end{proof}

\begin{lemma}
\label{lem:nu-h2^2gA'}
\revdeg{87, 12, 48}
There is a hidden $\nu$ extension from
$h_2^2 g A'$ to $\D h_1^2 g_2 g$.
\end{lemma}

\begin{proof}
Comparison to $C\tau$ shows that $h_2^2 g A'$ supports a hidden
$\nu$ extension whose target is either $\D h_1^2 g_2 g$
or $\D h_1^2 g_2 g + \tau h_2 g C''$.

Let $\alpha$ be an element of $\pi_{84,46}$ that is detected
by $h_2 g A'$.  Since $h_2 g A'$ does not support a hidden
$\eta$ extension, we may choose $\alpha$ such that $\eta \alpha$ is
zero.  Note that $h_2^2 g A'$ detects $\nu \alpha$.

Shuffle to obtain
\[
\nu^2 \alpha = \langle \eta, \nu, \eta \rangle \alpha =
\eta \langle \nu, \eta, \alpha \rangle.
\]
This shows that $\nu^2 \alpha$ must be divisible by $\eta$.
Consequently, the hidden $\nu$ extension on 
$h_2^2 g A'$ must have target $\D h_1^2 g_2 g$.
\end{proof}

\begin{remark}
\label{rem:nu,eta,h2gA'}
The proof of Lemma \ref{lem:nu-h2^2gA'} shows that
$\D h_1 g_2 g$ detects the Toda bracket
$\langle \nu, \eta, \{h_2 g A'\} \rangle$.
\end{remark}

\section{Miscellaneous hidden extensions}
\label{sctn:misc-extn}

\rev{
\begin{thm}
\label{thm:misc-extn}
Tables \ref{tab:misc-extn} \rev{and \ref{tab:misc-extn-null}} list some miscellaneous hidden extensions.
\end{thm}

\begin{proof}
Similarly to Theorems \ref{thm:2-extn}, \ref{thm:eta-extn}, and \ref{thm:nu-extn}, some of the extensions follow by comparison to $C\tau$ or to $\tmf$.
The more difficult cases are handled in the following lemmas.
\end{proof}
}

\rev{
Based on the corrected statement of Lemma~\ref{lem:theta4,2,sigma^2+kappa} (\cite{Xu16}*{Theorem 2.1}) and Lemma~\ref{lem:etakappabar2,2sigma,sigma}, the third author presents the proof of the following Lemma~\ref{lem:theta4-h4^2} (\cite{Xu16}*{Theorem 1.2}), fixing a gap in its original proof. Note that as in the original proof, it only uses classical knowledge back then up to the 60-stem.

\begin{lemma}
\label{lem:theta4-h4^2}
\revdeg{30, 2}
Classically, there is no hidden $\theta_4$ extension on $h_4^2$.
In other words $\theta_4^2$ is zero in $\pi_{60}$.
\end{lemma}

\begin{proof}
We have
$$\theta_4^2 = \theta_4 \langle 2, \sigma^2 + \kappa, 2\sigma, \sigma \rangle \subseteq \langle \langle \theta_4, 2, \sigma^2 + \kappa \rangle, 2\sigma, \sigma \rangle.$$
Using \cite{Xu16}*{Theorem~2.2} and Lemma \ref{lem:theta4,2,sigma^2+kappa}, the last expression is contained in the union of 
$$\langle 0, 2\sigma, \sigma \rangle, \ \langle \eta\kappabar_2, 2\sigma, \sigma \rangle, \ \langle \rho_{15}\theta_4, 2\sigma, \sigma \rangle.$$
By \cite{Xu16}*{Lemmas~2.3 and 2.4} and Lemma~\ref{lem:etakappabar2,2sigma,sigma}, all three brackets contain a single element 0. Therefore, $\theta_4^2 = 0$.
\end{proof}
}

\begin{lemma}
\label{lem:epsilon-h3^2h5}
\mbox{}
\begin{enumerate}
\item
\revdeg{45, 3, 24}
There is a hidden $\epsilon$ extension from $h_3^2 h_5$ to $M c_0$.
\item
\revdeg{45, 3, 23}
There is a hidden $\epsilon$ extension from $\tau h_3^2 h_5$ to $M P$.
\end{enumerate}
\end{lemma}

\begin{proof}
Table \ref{tab:eta-extn} shows that
$M h_1$ detects the product $\eta \theta_{4.5}$.
Then $M h_1 c_0$ detects $\eta \epsilon \theta_{4.5}$.
This implies that $M c_0$ detects $\epsilon \theta_{4.5}$.

This only shows that $M c_0$ is the target of a hidden
$\epsilon$ extension, whose source could be $h_3^2 h_5$
or $h_5 d_0$.  However, Lemma \ref{lem:epsilon-h5d0}
rules out the latter case.
This establishes the first hidden extension.

Table \ref{tab:tau-extn} shows that there is a hidden
$\tau$ extension from $M c_0$ to $M P$.  Then the first
hidden extension implies the second one.
\end{proof}

\begin{remark}
We claimed in \cite{Isaksen14c}*{Table 33} that there is a
hidden $\epsilon$ extension from $h_3^2 h_5$ to $M c_0$.
However, the argument given in \cite{Isaksen14c}*{Lemma 4.108}
only implies that $M c_0$ is the target of a hidden extension
from either $h_3^2 h_5$ or $h_5 d_0$.
\end{remark}

\begin{lemma}
\label{lem:kappa-h3^2h5}
\revdeg{45, 3, 24}
There is a hidden $\kappa$ extension from $h_3^2 h_5$ to $M d_0$.
\end{lemma}

\begin{proof}
Table \ref{tab:eta-extn} shows that
$M h_1$ detects the product $\eta \theta_{4.5}$.
Then $M h_1 d_0$ detects the product $\eta \kappa \theta_{4.5}$,
so $M d_0$ must detect the product $\kappa \theta_{4.5}$.
This shows that $M d_0$ is the target of a hidden $\kappa$ extension
whose source is either $h_3^2 h_5$ or $h_5 d_0$.

We showed in Lemma \ref{lem:epsilon-h5d0}
that $\epsilon \alpha$ is zero for some element 
$\alpha$ of $\pi_{45,24}$ that is detected by $h_5 d_0$.
Then $\epsilon \kappabar \alpha$ is also zero.
Table \ref{tab:misc-extn} shows that
$\epsilon \kappabar$ equals $\kappa^2$.
Therefore,
$\kappa^2 \alpha$ is zero.
If $\kappa \alpha$ were detected by $M d_0$, 
then $\kappa^2 \alpha$ would be detected by $M d_0^2$.
It follows that there is no hidden $\kappa$ extension
from $h_5 d_0$ to $M d_0$.
\end{proof}

\begin{remark}
We showed in \cite{Isaksen14c}*{Table 33} that there is a 
hidden $\kappa$ extension from either $h_3^2 h_5$ or $h_5 d_0$
to $M d_0$. Lemma \ref{lem:kappa-h3^2h5} settles this uncertainty.
\end{remark}

\begin{lemma}
\label{lem:kappabar-h3^2h5}
\revdeg{45, 3, 24}
There is a hidden $\kappabar$ extension from $h_3^2 h_5$ to
$\tau M g$.
\end{lemma}

\begin{proof}
Table \ref{tab:eta-extn} shows that
$M h_1$ detects the product $\eta \theta_{4.5}$.
Then $\tau M h_1 g$ detects the product $\eta \kappabar \theta_{4.5}$,
so $\tau M g$ must detect the product $\kappabar \theta_{4.5}$.
This shows that $\tau M g$ is the target of a hidden 
$\kappabar$ extension
whose source is either $h_3^2 h_5$ or $h_5 d_0$.

We showed in Lemma \ref{lem:epsilon-h5d0}
that $\epsilon \alpha$ is zero for some element 
$\alpha$ of $\pi_{45,24}$ that is detected by $h_5 d_0$.
If $\kappabar \alpha$ were detected by $\tau M g$,
then $\epsilon \kappabar \alpha$ would be detected by
$M d_0^2$ because Table \ref{tab:misc-extn} shows that
there is a hidden $\epsilon$ extension from $\tau M g$
to $M d_0^2$.
Therefore, there is no hidden $\kappabar$ extension
from $h_5 d_0$ to $\tau M g$.
\end{proof}

\begin{lemma}
\label{lem:Dh1h3-h3^2h5}
\revdeg{45, 3, 24}
There is a hidden $\{\D h_1 h_3\}$ extension from
$h_3^2 h_5$ to $M \D h_1 h_3$.
\end{lemma}

\begin{proof}
Table \ref{tab:eta-extn} shows that $M h_1$ detects the product
$\eta \theta_{4.5}$.  Therefore, the element
$M \D h_1^2 h_3$ detects $\eta \{\D h_1 h_3\} \theta_{4.5}$.
This shows that $M \D h_1 h_3$ is the target of a hidden
$\{\D h_1 h_3\}$ extension.
Lemma \ref{lem:Dh1h3-h5d0} rules out $h_5 d_0$ as a possible source.
The only remaining possible source is $h_3^2 h_5$.
\end{proof}

\begin{lemma}
\label{lem:theta4.5-h3^2h5}
\revdeg{45, 3, 24}
There is a hidden $\theta_{4.5}$ extension from
$h_3^2 h_5$ to $M^2$.
\end{lemma}

\begin{proof}
The proof of Lemma \ref{lem:perm-M^2} shows that
$M^2 h_1$ detects a multiple of $\eta \theta_{4.5}$.
Therefore, it detects either $\eta \theta_{4.5}^2$
or $\eta \theta_{4.5} \{h_5 d_0\}$.

Now $h_1 h_5 d_0$ detects $\eta \{h_5 d_0\}$, which also
detects $\eta_4 \theta_4$ by Table \ref{tab:misc-extn}.
In fact, the proof of \cite{Isaksen14c}*{Lemma 4.112} shows that
these two products are equal.
Then $\eta \theta_{4.5} \{h_5 d_0\}$ equals
$\eta_4 \theta_4 \theta_{4.5}$.
Next,
$\eta_4 \theta_{4.5}$ lies in $\pi_{61,33}$.
The only non-zero element of $\pi_{61,33}$ is detected by
$\mmf$, so the product $\eta_4 \theta_{4.5}$ must be zero.

We have now shown that $M^2 h_1$ detects $\eta \theta_{4.5}^2$.
This implies that $M^2$ detects $\theta_{4.5}^2$.
\end{proof}

\begin{lemma}
\label{lem:epsilon-h5d0}
\revdeg{45, 5, 24}
There is no hidden $\epsilon$ extension on $h_5 d_0$.
\end{lemma}

\begin{proof}
Table \ref{tab:Toda} shows that
$h_5 d_0$ detects the Toda bracket $\langle 2, \theta_4, \kappa \rangle$.
Now shuffle to obtain
\[
\epsilon \langle 2, \theta_4, \kappa \rangle =
\langle \epsilon, 2, \theta_4 \rangle \kappa.
\]
Table \ref{tab:Toda} shows that
$h_5 c_0$ detects the Toda bracket
$\langle \epsilon, 2, \theta_4 \rangle$, and there is no indeterminacy.
Let $\alpha$ in $\pi_{39,21}$ be the unique element of this Toda bracket.
We wish to compute $\alpha \kappa$.

Table \ref{tab:Toda} shows that
$h_5 c_0$ also detects the Toda bracket 
$\langle \eta_5, \nu, 2 \nu \rangle$, 
with indeterminacy generated by $\sigma \eta_5$.
Let $\beta$ in $\pi_{39,21}$ be an element of
this Toda bracket.
Then $\alpha$ and $\beta$ are equal, modulo
$\sigma \eta_5$ and modulo elements in higher filtration.
Both $\tau h_3 d_1$ and $\tau^2 c_1 g$ detect multiples of
$\sigma$.
Also, the difference between $\alpha$ and $\beta$ cannot be
detected by $\D h_1 d_0$ by comparison to $\tmf$.

This implies that $\alpha$ equals $\beta + \sigma \gamma$
for some element $\gamma$ in $\pi_{31,17}$.
Then
\[
\alpha \kappa = (\beta + \sigma \gamma) \kappa = \beta \kappa
\]
because $\sigma \kappa$ is zero.

Now shuffle to obtain
\[
\beta \kappa = \langle \eta_5, \nu, 2 \nu \rangle \kappa =
\eta_5 \langle \nu, 2\nu, \kappa \rangle.
\]
Table \ref{tab:Toda} shows that
$\langle \nu, 2\nu, \kappa \rangle$
contains $\eta \kappabar$, and its indeterminacy is generated by
$\nu \nu_4$.
We now need to compute $\eta_5 \eta \kappabar$.

The product $\eta_5 \kappabar$ is detected by
$\tau h_1 h_3 g_2 = \tau h_1 h_5 g$, so
$\eta_5 \kappabar$ equals $\tau \eta \sigma \kappabar_2$,
modulo elements of higher filtration.
But these elements of higher filtration are either annihilated
by $\eta$ or detected by $\tmf$, so 
$\eta_5 \eta \kappabar$ equals
$\tau \eta^2 \sigma \kappabar_2$.
By comparison to $\tmf$, this latter expression must be zero.
\end{proof}

\begin{lemma}
\label{lem:Dh1h3-h5d0}
\revdeg{45, 5, 24}
There is no hidden $\{\D h_1 h_3\}$ extension on 
$h_5 d_0$.
\end{lemma}

\begin{proof}
Table \ref{tab:Toda} shows that $h_5 d_0$ detects the Toda bracket
$\langle \kappa, \theta_4, 2 \rangle$.  By inspection of 
indeterminacies, we have
\[
\{\D h_1 h_3\} \langle \kappa, \theta_4, 2 \rangle =
\langle \{\D h_1 h_3\} \kappa, \theta_4, 2 \rangle.
\]
Table \ref{tab:misc-extn} shows that
$\tau d_0 l + \D c_0 d_0$ detects the product
$\{\D h_1 h_3 \} \kappa$.
Now apply the Moss Convergence Theorem \ref{thm:Moss}
with the Adams differential $d_2(h_5) = h_0 h_4^2$
to determine that the Toda bracket
$\langle \{\D h_1 h_3\} \kappa, \theta_4, 2 \rangle$ is 
detected in Adams filtration at least $13$.

The only element in sufficiently high filtration is $\tau^5 e_0 g^3$,
but comparison to $\mmf$ rules this out.
Thus the Toda bracket
$\langle \{\D h_1 h_3\} \kappa, \theta_4, 2 \rangle$ 
contains zero.
\end{proof}

\begin{lemma}
\label{lem:rho15-h5^2}
\revdeg{62, 2, 32}
There is a hidden $\rho_{15}$ extension from $h_5^2$ to
either $h_0 x_{77,7}$ or $\tau^2 m_1$.
\end{lemma}

\begin{proof}
Table \ref{tab:Toda} shows that the Toda bracket
$\langle 8, 2 \sigma, \sigma \rangle$ contains
$\rho_{15}$.
Then $\rho_{15}$ is also contained in 
$\langle 2, 8 \sigma, \sigma \rangle$, although the indeterminacy
increases.

Now shuffle to obtain
\[
\rho_{15} \theta_5 = \theta_5 \langle 2, 8\sigma, \sigma \rangle =
\langle \theta_5, 2, 8 \sigma \rangle \sigma.
\]
Table \ref{tab:Toda} shows that $h_0^3 h_3 h_6$
detects $\langle \theta_5, 2, 8\sigma \rangle$.
Also, there is a $\sigma$ extension from $h_0^3 h_3 h_6$
to $h_0 x_{77,7}$ in the homotopy of $C\tau$.

This implies that $\rho_{15} \theta_5$ is non-zero in
$\pi_{77,40}$, and that it is detected in filtration at most $8$.
Moreover, it is detected in filtration at least $7$, since
$\rho_{15}$ and $\theta_5$ are detected in filtrations $4$ and $2$
respectively.

There are several elements
in filtration $7$ that could detect $\rho_{15} \theta_5$.
The element $x_{77,7}$ (if it survives to the $E_\infty$-page)
is ruled out by comparison to $C\tau$.
The element $\tau h_1 x_{76,6}$ is ruled out because
$\eta \rho_{15} \theta_5$ is detected in filtration at least $10$,
since $\eta \rho_{15}$ is detected in filtration $7$.

The only remaining possibilities are $h_0 x_{77,7}$
and $\tau^2 m_1$.
\end{proof}

\begin{lemma}
\label{lem:rho15-h1h6}
\mbox{}
\begin{enumerate}
\item
\revdeg{64, 2, 33}
There is a hidden $\rho_{15}$ extension from
$h_1 h_6$ to $P h_6 c_0$.
\item
\revdeg{64, 2, 33}
There is a hidden $\rho_{23}$ extension from
$h_1 h_6$ to $P^2 h_6 c_0$.
\end{enumerate}
\end{lemma}

\begin{proof}
Table \ref{tab:Toda} shows that
$h_1 h_6$ detects $\langle \eta, 2, \theta_5 \rangle$.
Then
\[
\rho_{15} \langle \eta, 2, \theta_5 \rangle \subseteq
\langle \eta \rho_{15}, 2, \theta_5 \rangle.
\]
The last bracket is detected by $P h_6 c_0$ because
$d_2(h_6) = h_0 h_5^2$ and because
$P c_0$ detects $\eta \rho_{15}$.
Also, its indeterminacy is in Adams filtration greater than $8$.
This establishes the first hidden extension.

The proof for the second extension is essentially the same,
using that 
$P^2 c_0$ detects $\eta \rho_{23}$ and that
the indeterminacy of
$\langle \eta \rho_{23}, 2, \theta_5 \rangle$
is in Adams filtration greater than $12$.
\end{proof}


\begin{lemma}
\label{lem:epsilon-tMg}
\revdeg{65, 10, 35}
There is a hidden $\epsilon$ extension from $\tau M g$ to
$M d_0^2$.
\end{lemma}

\begin{proof}
First, we have the relation $c_0 \cdot h_1^2 X_2 = M h_1 h_3 g$
in the Adams $E_2$-page, which is detected in the homotopy of
$C\tau$.
Table \ref{tab:tau-extn} shows that there are hidden $\tau$
extensions from $h_1^2 X_2$ and $M h_1 h_3 g$
to $\tau M g$ and $M d_0^2$ respectively.
\end{proof}

\begin{lemma}
\label{lem:epsilon-MDh1h3}
\revdeg{77, 12, 41}
If $M \D h_1^2 d_0$ is non-zero in the $E_\infty$-page, then
there is a hidden $\epsilon$ extension from $M \D h_1 h_3$
to $M \D h_1^2 d_0$.
\end{lemma}

\begin{proof}
Table \ref{tab:eta-extn} shows that $M h_1$ detects the 
product $\eta \theta_{4.5}$.
Since $M \D h_1^2 d_0$ equals $\D h_1 d_0 \cdot M h_1$, it detects
$\eta \{\D h_1 d_0\} \theta_{4.5}$.

Table \ref{tab:misc-extn} shows that 
$\eta \{\D h_1 d_0\}$ equals 
$\epsilon \{\D h_1 h_3\}$, since they are both
detected by $\D h_1^2 d_0$ and there are no elements in higher
Adams filtration.
Therefore, the product 
$\epsilon \{\D h_1 h_3\} \theta_{4.5}$ is detected
by $M \D h_1^2 d_0$.
In particular, $\{\D h_1 h_3\} \theta_{4.5}$ is non-zero, and
it can only be detected by $M \D h_1 h_3$.
\end{proof}

\section{Additional relations}

\begin{lemma}
\label{lem:nubar-h5^2}
The product $(\eta \sigma + \epsilon) \theta_5$ is detected by
$\tau h_0 h_2 Q_3$.
\end{lemma}

\begin{proof}
Table \ref{tab:Toda} indicates a hidden
$\eta$ extension from $h_2^2 h_6$ to $\tau h_0 h_2 Q_3$.
Therefore, there exists an element $\alpha$ in $\pi_{69,36}$
such that $\tau h_0 h_2 Q_3$ detects $\eta \alpha$.
(Beware of the crossing extension from $p'$ to $h_1 p'$.
This means that it is possible to choose such an $\alpha$, but not
any element detected by $h_2^2 h_6$ will suffice.)

Table \ref{tab:Toda} shows that $h_2^2 h_6$ detects the
Toda bracket $\langle \nu^2, 2, \theta_5 \rangle$.
Let $\beta$ be an element of this Toda bracket.
Since $\alpha$ and $\beta$ are both detected by $h_2^2 h_6$, 
the difference $\alpha - \beta$ is
detected in Adams filtration at least $4$.

Table \ref{tab:misc-extn} shows that
$p'$ detects $\sigma \theta_5$, which belongs to the 
indeterminacy of $\langle \nu^2, 2, \theta_5 \rangle$.
Therefore, we may choose $\beta$ such that
the difference $\alpha - \beta$
is detected in filtration at least 9.
Since $\eta \alpha$ is detected by $\tau h_0 h_2 Q_3$ in filtration 7,
it follows that $\eta \beta$ is also detected by $\tau h_0 h_2 Q_3$.

We now have an element $\beta$ contained in
$\langle \nu^2, 2, \theta_5 \rangle$ such that
$\eta \beta$ is detected by $\tau h_0 h_2 Q_3$.
Now consider the shuffle
\[
\eta \langle \nu^2, 2, \theta_5 \rangle =
\langle \eta, \nu^2, 2 \rangle \theta_5.
\]
Table \ref{tab:Toda} shows that the last bracket equals
$\{ \epsilon, \epsilon + \eta \sigma \}$.
Therefore, either
$\epsilon \theta_5$ or
$(\epsilon + \eta \sigma) \theta_5$ is detected by
$\tau h_0 h_2 Q_3$.
But $\epsilon \theta_5$ is detected by $h_1 p' = h_5^2 c_0$.
\end{proof}

\begin{lemma}
\label{lem:compound-(epsilon+etasigma)-theta5}
There exists an element $\alpha$ of $\pi_{67,36}$ that is 
detected by $h_0 Q_3 + h_0 n_1$ such that
$\tau \nu \alpha$ equals $(\eta \sigma + \epsilon) \theta_5$.
\end{lemma}

\begin{proof}
Lemma \ref{lem:nubar-h5^2} shows that
$\tau h_0 h_2 Q_3$ detects $(\epsilon + \eta \sigma) \theta_5$.
The element $\tau h_0 h_2 Q_3$ also detects $\tau \nu \alpha$.
Let $\beta$ be the difference
$\tau \nu \alpha - (\epsilon + \eta \sigma) \theta_5$,
which is detected in higher Adams filtration.
We will show that $\beta$ must equal zero.

First, $\tau h_1 D'_3$ cannot detect $\beta$ because
$\eta^2 \beta$ is zero, while 
Table \ref{tab:eta-extn} shows that $\tau^3 d_1 g^2$ detects
$\eta^2 \{\tau h_1 D'_3\}$.
Second, Table \ref{tab:nu-extn} shows that $\tau h_1 h_3(\D e_1 + C_0)$
is the target of a hidden $\nu$ extension.
Therefore, we may alter the choice of $\alpha$
to ensure that $\beta$ is not detected by $\tau h_1 h_3(\D e_1 + C_0)$.
Third, $\D^2 h_2 c_1$ is also the target of a $\nu$ extension.
Therefore, we may alter the choice of $\alpha$
to ensure that $\beta$ is not detected by $\D^2 h_2 c_1$.
Finally, comparison to $\tmf$ implies that
$\beta$ is not detected by 
$\tau \D^2 h_1^2 g + \tau^3 \D h_2^2 g^2$. 
\end{proof}



\rev{
\begin{lemma}
\label{lem:sigma^2+kappa-h5^2}
The product $(\sigma^2 + \kappa) \theta_5$ is
zero, or it is equal to $\tau^2 \kappa_1 \kappabar_2$ detected
by $\tau^2 d_1 g_2$.
\end{lemma}



\begin{proof}
The product
$(\sigma^2 + \kappa) \theta_5$ maps to zero under
inclusion of the bottom cell of $C\tau$.
Therefore, $(\sigma^2 + \kappa) \theta_5$ is divisible
by $\tau$.  The only two possibilities are $0$ and
$\tau^2 \kappa_1 \kappabar_2$.
\end{proof}
}

\begin{lemma}
\label{lem:eta-sigma-k1}
$\eta \sigma \{k_1\} + \nu \{d_1 e_1\} = 0$ in
$\pi_{73,41}$.
\end{lemma}

\begin{proof}
We have the relation $h_1 h_3 k_1 + h_2 d_1 e_1 = 0$
in the Adams $E_\infty$-page, but
$\eta \sigma \{k_1\} + \nu\{d_1 e_1\}$ could possibly be 
detected in higher Adams filtration.
However, it cannot be detected by
$h_2^2 C''$ or $M h_1 h_3 g$ by comparison to $C\tau$.
Also, it cannot be detected by $\D h_1 d_0 e_0^2$ by comparison
to $\mmf$. 
\end{proof}
